% !TeX program = xelatex
% This report's numerical tables and chart data are generated by Dask.
% The primary analysis and historical evidence remain separate source files.
% Shared XeLaTeX report style. Charts consume Dask-produced CSVs via PGFPlots.
\documentclass[a4paper,fontsize=11pt,parskip=half,headings=small]{scrartcl}
\usepackage[margin=18mm,includeheadfoot,headheight=14pt,headsep=6mm,footskip=9mm]{geometry}
\usepackage{fontspec}
\usepackage{mathtools,amsthm}
\usepackage{unicode-math}
\setmainfont{TeX Gyre Pagella}
\setsansfont{TeX Gyre Heros}[Scale=MatchLowercase]
\setmonofont{Latin Modern Mono}
\setmathfont{TeX Gyre Pagella Math}
\usepackage[protrusion=true,expansion=false]{microtype}
\usepackage{xcolor,xurl,booktabs,longtable,array,needspace,placeins}
\usepackage[numbers,sort&compress]{natbib}
\usepackage{pgfplots}
\pgfplotsset{compat=1.18,width=.92\linewidth,height=6.0cm,
 every axis/.append style={grid=major,grid style={gray!20},
 tick label style={font=\small},label style={font=\small},
 legend style={font=\small,draw=none,fill=white}}}
\usepgfplotslibrary{groupplots}
\usepackage[unicode,colorlinks=true,linkcolor=blue!50!black,citecolor=blue!50!black,urlcolor=blue!50!black]{hyperref}
\usepackage{bookmark,scrlayer-scrpage}
\setlength{\emergencystretch}{1.5em}
\setkomafont{disposition}{\normalfont\sffamily\bfseries}
\setkomafont{pageheadfoot}{\normalfont\sffamily\small}
\clearpairofpagestyles
\ihead{Simulation methods and results}
\ohead{Dask experiment report}
\cfoot*{\pagemark}
\pagestyle{scrheadings}
\newtheorem{theorem}{Theorem}[section]
\newtheorem{proposition}[theorem]{Proposition}
\newtheorem{lemma}[theorem]{Lemma}
\newtheorem{corollary}[theorem]{Corollary}
\DeclareMathOperator{\E}{E}
\DeclareMathOperator{\Var}{Var}
\DeclareMathOperator{\Cov}{Cov}
\DeclareMathOperator{\SD}{SD}

\title{Hoeflin/Ferguson: Simulation Methods and Results}
\subtitle{SPARK and Dask experiments: composites, prediction and rare tails}
\author{\ProjectAuthor}
\date{16 November 2026}

\begin{document}
\frontmatter
\maketitle
\thispagestyle{plain.scrheadings}
% Computed by the Dask report task; display rounding only.
\newcommand{\SimTrials}{1\,000\,000}
\newcommand{\SimReplicates}{4}
\newcommand{\SimPooledTrials}{4\,000\,000}
\newcommand{\SimChunk}{50\,000}
\newcommand{\SimMaxMeanSE}{0.661}
\newcommand{\SimMaxVarianceRatio}{0.000254}
\newcommand{\SimRareMin}{1}
\newcommand{\SimRareMax}{4}
\newcommand{\SimMaxEventSE}{2.308}
\newcommand{\SimMaxSDRelativeError}{0.000127}


\begin{abstract}
This report explains both the SPARK/Ada simulations accompanying the
Hoeflin/Ferguson analysis and their Dask reproduction under the same
expressly assumed Gaussian common-factor model. The source-pinned native
experiment supplies one million observations for each of eight battery
sizes at correlation \(0.6\). In Dask, four seeded replicates supply \SimTrials\ observations
per size per replicate, or \SimPooledTrials\ observations for each fixed
size. We examine the variance of the unstandardized mean, the calibration
of the standardized composite, prediction of the simulated common factor,
and exceedance counts at moderate and one-in-a-million thresholds. The
numerical results, algorithms, estimators and verification boundaries are
examined for both implementations against exact Gaussian targets. The
native experiment reports composite moments and tail counts; Dask also
measures latent-factor prediction error. Dask constructs and executes the numerical task graph;
XeLaTeX/PGFPlots renders charts from its output tables. No real test data
are generated or analysed, and numerical agreement does not establish the
psychometric model.
\end{abstract}

\noindent\textbf{AI disclosure.} This exposition, its supporting code and its
review were prepared with OpenAI Codex and collaborating AI agents. The
requester Anders Hellström specified the report and verification requirements. The work does
not claim independent human peer review or empirical validation of IQ tests.
The full disclosure appears in \cref{sec:ai-disclosure}.

\noindent\textbf{Proof and evidence convention.} The model deductions in
this report are proved in \cref{sec:proof-foundations,hfpr:appendix};
\Cref{hfmap:simulation} maps every claim to
its proof, model hypothesis or recorded evidence. Numerical reference
values and observed estimates are computed approximations. Neither a
matching decimal nor a passing program contract proves human-score
Gaussianity or a certified last-digit numerical bound.

\tableofcontents
\mainmatter

\chapter{Questions, experimental units and evidence}
The main analysis revisits \citet{hoeflin}'s use of covariance identities to
combine test scores. Its simulation asks what happens \emph{inside a chosen
probability model}; it does not estimate a population correlation from the
27 selected scorers in the historical article. We fix every test's
population mean at zero, its SD at one, and every distinct pairwise
correlation at \(c=0.6\). The simulated quantities answer three questions:

\begin{enumerate}
\item Does the raw test-score average have the variance predicted by the
covariance calculation, and does dividing by its population SD produce a
composite with mean zero and variance one?
\item In the particular one-factor Gaussian model, how accurately does the
conditional-mean predictor recover the \emph{simulated} latent factor?
\item How much information do the observed exceedance counts contain about
moderate tails and an assumed upper-tail probability of \(10^{-6}\)?
\end{enumerate}

An observation is one synthetic draw of an aggregate test-score mean and its
common factor. It is not a human participant or a new empirical norming
observation. The eight sizes are
\[n\in\{1,2,3,4,5,10,100,1000\}.\]
For each fixed size the pooled sample contains
\(T=\SimPooledTrials\) observations from \SimReplicates\ seeded replicates.
The same standard-normal arrays are reused across battery sizes within a
replicate. Thus the eight reported results are coupled comparisons; they
are not eight independent replications of a population experiment.

Two sources of simulation evidence remain distinct. The Dask results use
the recorded random generator, seeds, chunks and software versions. The
SPARK/Ada experiment contains one million observations per size, with the
seed reset to \texttt{20261030} for every size. Its final 4 November native
runtime reproduced the original archived CSV byte for byte; this report
therefore discusses the full native experiment without counting that
repeat as fresh Monte Carlo information. Different generators and
execution plans give different individual draws. Comparison concerns the
defined distributions, estimators and sampling precision, rather than
equality of counts. \Cref{sec:hf-spark-study} documents the native methods,
results, proof scope and exact retained source context.

\chapter{Model and exact targets}
\section{From individual tests to the simulated aggregate}
Suppose
\begin{equation}
 X_i=\sqrt c\,G+\sqrt{1-c}\,\varepsilon_i,
 \qquad i=1,\ldots,n,
 \label{eq:hf-factor-model}
\end{equation}
where \(G,\varepsilon_1,\ldots,\varepsilon_n\) are independent standard
normals. Independence and the moment identities give
\(\operatorname E[X_i]=0\),
\(\operatorname{Var}(X_i)=c+(1-c)=1\), and
\(\operatorname{Cov}(X_i,X_j)=c\) for \(i\ne j\).
For \(M_n=n^{-1}\sum_iX_i\), covariance bilinearity therefore gives
\begin{equation}
 v_n=\operatorname{Var}(M_n)
 =\frac{n+n(n-1)c}{n^2}
 =c+\frac{1-c}{n}.
 \label{eq:hf-average-variance}
\end{equation}
These moment calculations require finite second moments and the stated
covariances. Normality is used for the stronger distributional statements
below. They follow from Gaussian affine closure
\citep[Definition 3.3.1 and Theorem 3.3.1]{gallager2013}.
Propositions~\ref{hfpr:standardization}, \ref{hfpr:equal-weights} and
\ref{hfpr:factor} provide the full moment and distributional proofs.

The Dask simulation uses the exact aggregate representation
\begin{equation}
 M_n\ \stackrel{d}{=}\ \sqrt c\,G+
        \sqrt{\frac{1-c}{n}}\,E,
 \qquad G,E\ \hbox{independent standard normals}.
 \label{eq:hf-aggregate}
\end{equation}
Indeed, \(n^{-1/2}\sum_i\varepsilon_i\) is standard normal and independent
of \(G\), proving equality in distribution with
\eqref{eq:hf-factor-model}'s mean. Both constructions also have the same
joint law for \((G,M_n)\), since their cross-covariance is \(\sqrt c\).
Consequently the shortcut preserves the targets used in this report; it
avoids generating a thousand individual residual scores for every draw of
the thousand-test case.
The representation is proved jointly with \(G\) in
Proposition~\ref{hfpr:factor}.

Reusing one \(E\) across different sizes is a computational coupling of
these correct \emph{fixed-size} laws. It does not assert that the simulated
paths are the successive averages of a single nested battery of actual
individual scores. That distinction matters when interpreting correlations
between the plotted points or adding evidence across sizes.
Proposition~\ref{hfpr:simulation} derives the different cross-size
covariances in the coupled shortcut and a nested battery explicitly.

\section{Standardized composites and qualifying thresholds}
Define
\begin{equation}
 Z_{C,n}=\frac{M_n}{\sqrt{v_n}}.
 \label{eq:hf-composite}
\end{equation}
Equations~\eqref{eq:hf-average-variance} and \eqref{eq:hf-composite}
imply mean zero and variance one. Gaussian affine closure additionally
makes \(Z_{C,n}\) standard normal. Therefore every size has the same exact
target exceedance probabilities:
\begin{align}
 P(Z_{C,n}>2)&=\overline\Phi(2),&
 P(Z_{C,n}>3)&=\overline\Phi(3),\\
 P(Z_{C,n}>q)&=10^{-6},&
 q&=\overline\Phi^{-1}(10^{-6}).
 \label{eq:hf-tail-targets}
\end{align}
The evaluator uses SciPy's upper-tail inverse directly, avoiding subtraction
of a probability close to one \citep{scipystats}.

The corresponding cutoff on the raw average is \(q\sqrt{v_n}\), since
\(M_n>q\sqrt{v_n}\) is exactly the event \(Z_{C,n}>q\).
Thus the raw cutoff falls with \(n\), while the population percentile
assigned by this assumed model remains fixed. An IQ-style reporting scale
\(100+16M_n\) changes units; it does not make that average a unit-SD
composite. Conversely, the standardized composite's SD-16 reporting value
at the rare cutoff is always \(100+16q\). These follow from affine
expectation and variance identities \citep[Section 1.6]{durrett2019}.

At fixed \(c=0.6\), \(v_n\) tends to \(c\), because
\(v_n-c=(1-c)/n\). Its square root tends to \(\sqrt c\), by continuity
of the nonnegative square root. Raw average SD therefore has a positive
limiting floor, while composite SD remains one. This is a consequence of
the common uncertain component, rather than a cap on possible scores.
Propositions~\ref{hfpr:thresholds} and \ref{hfpr:numerics} prove the
event equivalence, monotonicity, limiting error bound and scale conversions.

\section{Prediction of the simulated common factor}
In this model, \(\operatorname{Cov}(G,M_n)=\sqrt c\). Put
\begin{equation}
 \widehat G_n=\frac{\sqrt c}{v_n}M_n,
 \qquad R_n=\frac{c}{v_n}
          =\frac{nc}{1+(n-1)c}.
 \label{eq:hf-predictor}
\end{equation}
The residual \(W_n=G-\widehat G_n\) is jointly Gaussian with \(M_n\)
and has covariance
\(\sqrt c-(\sqrt c/v_n)v_n=0\) with it. Gaussian uncorrelated blocks
are independent, so \(\operatorname E[G\mid M_n]=\widehat G_n\), and
\begin{equation}
 \operatorname E[(G-\widehat G_n)^2]
 =\operatorname{Var}(W_n)
 =1-\frac{c}{v_n}
 =\frac{1-c}{nc+1-c}.
 \label{eq:hf-mse}
\end{equation}
This is the Gaussian regression/residual argument
\citep[Theorems 3.5.1--3.5.2]{gallager2013}. The score contrasts
\(X_i-M_n\) have zero covariance with both \(G\) and \(M_n\), and
hence with \(W_n\); joint Gaussianity makes the same predictor valid
when conditioning on the full score vector. No extra information about
\(G\) is lost by retaining only the average in this symmetric model.

A standardized composite and a conditional factor predictor are distinct:
\(\widehat G_n=\sqrt{R_n}\,Z_{C,n}\), with
\(\operatorname{Var}(\widehat G_n)=R_n\), whereas
\(\operatorname{Var}(Z_{C,n})=1\). The conditional residual variance is
\(1-R_n\). As \(n\) increases, the latter tends to zero under this
particular independent-residual model, even though raw mean SD tends to
\(\sqrt c\). Shared population variation and prediction error answer
different questions. Nothing in this illustration establishes a factor
score correction for an actual IQ battery.
Proposition~\ref{hfpr:factor} proves the posterior law, the contrast
argument, predictor variance and MSE optimality from Gaussian independence.
Proposition~\ref{hfpr:nonidentification} proves why these conclusions
cannot identify an actual latent target from the score law alone.

\chapter{SPARK simulation: algorithm, results and verification}
\label{sec:hf-spark-study}
\section{Source context, cases and random-state construction}
The native evidence used here is the final Hoeflin run of the 4 November
2026 five-contract campaign. In repository-relative notation, its directory is
\begin{center}
\small\path{verification/evidence/five-contracts-computer10-20261104/}\par
\path{runtime-final/hoeflin/check/}
\end{center}
The relevant files are \path{simulation.csv}, \path{calculations.csv},
\path{runtime.json}, \path{validation-normal.json} and
\path{replay-normal.json}, with corresponding optimized-validator records.
The sibling \path{ignore/} directory records the Ghost Ignore execution.
The campaign's \path{manifest.json}, native portfolio result and
\path{runtime-final/summary.json} identify the source hashes, tools,
commands and policy comparison. The simulation CSV SHA-256 is
\begin{center}\small
\path{dba8b3046072a6f46b4c63117df9ac0b33aaa4f486507c493558391395c0ec36}.
\end{center}
It is byte-identical to the earlier archived CSV used as the Dask
campaign's historical input. No retained evidence file is rewritten by
this report revision, and no new native execution is claimed.

The SPARK driver evaluates all eight specified battery sizes at
\(c=0.6\), using \(T=1{,}000{,}000\) draws for each size. For every
size it resets a 64-bit state to seed \(20261030\). The recorded generator
advances by the modular recurrence
\begin{align*}
 s_{j+1}&=(A s_j+B)\bmod 2^{64},\\
 A&=6364136223846793005,\qquad B=1442695040888963407.
\end{align*}
It maps the high 52 bits into an open-unit value
\[
 U_j=\bigl(\lfloor s_j/2^{12}\rfloor+\tfrac12\bigr)2^{-52}.
\]
Because \(0\leq\lfloor s_j/2^{12}\rfloor\leq2^{52}-1\),
\(2^{-53}\leq U_j\leq1-2^{-53}\). This elementary bound explains
why the subsequent logarithm avoids zero; it does not establish a
continuous uniform law or independent successive values.

Two successive values \(U,V\) give the Box--Muller pair
\[
 G=\sqrt{-2\log U}\cos(2\pi V),\qquad
 E=\sqrt{-2\log U}\sin(2\pi V).
\]
The driver then computes \(M_n\) by \eqref{eq:hf-aggregate}, divides by
\(\sqrt{v_n}\), accumulates its first two sample moments, and updates the
three nested threshold counts. Resetting the seed reuses the same base
pairs across sizes. Consequently the SPARK experiment shares the same
fixed-size aggregate coupling described for Dask, rather than simulating
a nested sequence of independent individual-test residuals.
Proposition~\ref{hfpr:simulation} proves the coupling distinction.

\begin{proposition}[Ideal Box--Muller law]\label{hfsp:box-muller}
If \(U,V\) are independent continuous uniform variables on \((0,1)\),
the displayed Box--Muller formulas yield independent standard normals.
\end{proposition}
\begin{proof}
This is the classical transformation of \citet{boxmuller1958hf}; the
following change-of-variables argument gives the complete law. Put
\(R=\sqrt{-2\log U}\) and \(\Theta=2\pi V\). Since
\(U=\exp(-R^2/2)\), differentiation gives radial density
\(r\exp(-r^2/2)\) for \(r>0\). The angle has density \(1/(2\pi)\)
on \((0,2\pi)\), independently of \(R\). The polar map
\((r,\theta)\mapsto(g,e)=(r\cos\theta,r\sin\theta)\)
has absolute Jacobian \(r\), outside a measure-zero origin and ray.
The density change-of-variables theorem therefore gives joint density
\[
 \frac{1}{2\pi}\exp\!\left(-\frac{g^2+e^2}{2}\right)
 =\phi(g)\phi(e).
\]
Factorization into the two normalized marginal densities proves both
standard-normality and independence.
\end{proof}
The proposition is conditional on ideal continuous inputs. The finite
LCG, discrete open-unit mapping and external elementary functions are not
proved to satisfy those hypotheses. Indeed, the displayed lower bound
on \(U\) gives \(R\leq\sqrt{106\log2}\) in exact arithmetic and hence
\(|G|,|E|\leq\sqrt{106\log2}\), whereas a standard normal has positive
probability beyond every finite bound by its strictly positive density.
Thus the finite construction cannot literally have that continuous law. The theorem
justifies the intended simulation algorithm; it is separate from
implementation replay and range proofs.

\section{Estimators and the complete native results}
For the computed standardized samples \(Z_1,\ldots,Z_T\), the native
driver reports
\[
 \overline Z=\frac{S_1}{T},\qquad
 S_Z^2=\frac{S_2-S_1\overline Z}{T-1},\qquad
 S_1=\sum_j Z_j,\quad S_2=\sum_j Z_j^2,
\]
and \(K_b/T\), where \(K_b=\sum_j\mathbf1_{Z_j>b}\) for
\(b=2,3,q\). The variance numerator equals
\(\sum_j(Z_j-\overline Z)^2\), by expanding the square and using
\(\sum_jZ_j=T\overline Z\). Thus the estimator is the same unbiased
\(T-1\) sample variance used in the Dask pool. Unbiasedness and the
Gaussian mean/chi-square reference laws are proved in
Proposition~\ref{sm:normal-sample}; event moments are proved in
Proposition~\ref{sm:binomial}. Floating-point accumulation and printed
rounding remain implementation approximations to these exact formulas.

\begingroup\small

\begingroup\small
\setlength{\tabcolsep}{4pt}
\Needspace{6\baselineskip}
\noindent\textbf{Retained Ada run: one million observations per size}
\begin{longtable}{rrrrrr}
\toprule
$n$ & Archived mean & Archived variance & $K_2$ & $K_3$ & $K_q$ \\
\midrule
\endfirsthead
$n$ & Archived mean & Archived variance & $K_2$ & $K_3$ & $K_q$ \\
\midrule
\endhead
\bottomrule
\endfoot
$1$ & $3.969606\times10^{-4}$ & $0.999023$ & $22603$ & $1312$ & $1$ \\
$2$ & $5.006006\times10^{-4}$ & $0.999649$ & $22704$ & $1331$ & $0$ \\
$3$ & $5.491356\times10^{-4}$ & $0.999999$ & $22734$ & $1337$ & $1$ \\
$4$ & $5.780900\times10^{-4}$ & $1.000226$ & $22721$ & $1326$ & $1$ \\
$5$ & $5.976301\times10^{-4}$ & $1.000387$ & $22742$ & $1323$ & $1$ \\
$10$ & $6.444156\times10^{-4}$ & $1.000799$ & $22752$ & $1306$ & $1$ \\
$100$ & $7.115883\times10^{-4}$ & $1.001462$ & $22734$ & $1322$ & $1$ \\
$1000$ & $7.290209\times10^{-4}$ & $1.001649$ & $22685$ & $1325$ & $1$ \\
\end{longtable}
\endgroup

\endgroup
These eight rows are the full native simulation output, not selected
examples. Means range from approximately \(0.000397\) to \(0.000729\),
against mean target zero and ideal SE \(0.001\). Composite variances range
from \(0.999023\) to \(1.001649\), against target one. Their largest
absolute variance deviation is about \(0.001649\); the ideal variance
SE is approximately \(0.001414\). For completeness, the variance-SE
formula follows without an additional sampling assumption beyond
Proposition~\ref{sm:normal-sample}: write its chi-square variable as a sum
of \(T-1\) independent squared standard normals. For a standard normal
\(H\), integration by parts using \(\phi'(x)=-x\phi(x)\) gives
\(\E H^4=3\E H^2=3\); the boundary term
\(x^3\phi(x)\) vanishes at both infinities. Hence
\(\Var(H^2)=3-1=2\), and independence gives chi-square variance
\(2(T-1)\). Scaling by \((T-1)^{-2}\) yields
\(\Var(S_Z^2)=2/(T-1)\), whose square root is the stated SE.
The independent-normal representation is the referenced Gaussian
orthogonal-transformation proof in Proposition~\ref{sm:normal-sample}.
These finite-run deviations are retained
as observed, not rounded back to their theoretical values.

Moderate-tail counts range from \(22603\) to \(22752\) for \(Z>2\),
and from \(1306\) to \(1337\) for \(Z>3\). The respective expected
counts are approximately \(22750.132\) and \(1349.898\), obtained by
multiplying the targets in \eqref{eq:hf-tail-targets} by one million.
The native validator records 95\% Wilson intervals for all 16
moderate-tail proportions; all contain their corresponding Gaussian
reference. For example, the \(n=1\), \(Z>2\) interval is approximately
\([0.0223135,0.0228962]\), and its \(Z>3\) interval is
\([0.00124294,0.00138489]\). The score-inversion derivation and
approximate-coverage qualification are in Proposition~\ref{sm:wilson}.

The rare count is zero at \(n=2\) and one at every other size, against
expected count one for each fixed size. By the binomial variance formula,
its ideal count SD is \(\sqrt{1-10^{-6}}\), approximately one, and the
exact probability of no exceedances is \((1-10^{-6})^{10^6}\),
approximately \(e^{-1}\). Proposition~\ref{hfpr:simulation} proves the
zero-count identity and bounds the exponential approximation. The eight
coupled rows must not be pooled as eight million independent tail draws.

\section{Deterministic calculations accompanying the simulation}
\label{sec:hf-spark-calculations}
The companion \path{calculations.csv} supplies nine threshold rows:
the eight finite sizes and the limit \(n\to\infty\). Their reported
variance factors run from \(1\) at \(n=1\), through \(0.64\) at
\(n=10\), to \(0.6004\) at \(n=1000\), with limit \(0.6\).
The corresponding rare raw-mean cutoffs are approximately
\(4.753424309\), \(3.802739447\), \(3.683213761\), and
\(3.681986637\). These are calculations of \(q\sqrt{v_n}\), not
additional Monte Carlo estimates. Proposition~\ref{hfpr:thresholds}
proves the cutoff identity and limit.

Four further outputs check the rare quantile and worked examples. The
three-test mean of 160 on an SD-16 reporting scale gives composite value
\(170.064904975\). The latent conditional example reports mean
\(168.794145026\) and SD \(6.822422923\), on the separately specified
latent SD-16 scale. Proposition~\ref{hfpr:numerics} derives the exact
formulas and numerical interpretation. The normal and optimized native
validators report maximum absolute numerical error
\(2.842170943040401\times10^{-14}\), against tolerance
\(5\times10^{-11}\), over the nine rows and four outputs. This is
agreement with the validator's computed references; it is not a certified
last-digit bound on the external numerical libraries.

\section{Runtime replay, native proof and their boundaries}
The native runtime record reports successful GNAT 13.2.0 builds with
ordinary assertion and overflow checks enabled, both with Ghost Check
and Ghost Ignore. The two policies reproduce the same calculation and
simulation CSV bytes. Normal and optimized Python validators and the
separate-implementation replay pass; replay recovers all three tail
counts exactly in every size and matches means and unbiased variances
within its \(2\times10^{-12}\) tolerance. Eight CLI/domain cases are
also retained for each policy, including rejection of trial counts zero
and one. A repeated policy run or deterministic replay does not add new
random observations.

For the source hashes in that campaign, the native Hoeflin portfolio
discharged 404 checks: 129 by flow analysis and 275 by provers, with zero
unproved and zero justified checks. Native coverage records 53
\texttt{all} and 14 \texttt{spec} entities. The historical structural
inventory covers 35 functions and 28 procedures with 356 explicit
contract entries and a five-entry floor. These are different inventories:
entity/body classification, structural contract counting and discharged
proof obligations are not interchangeable measures. Nor does the
five-entry floor establish logical independence or adequacy of contracts.

The proved scope includes specified modular state advancement, bounded
arithmetic/accumulations, count updates, control flow and modeled report
receipts. Trusted numerical completions, \texttt{erfc}, logarithm,
square root, sine/cosine, compiler/runtime behavior and physical I/O retain
the boundaries documented in the campaign. Receipt contracts establish
modeled requests or normal returns, not successful physical output. No
proof here establishes generator independence, exact normal sampling,
inverse-tail accuracy or validity of the psychometric model.

The present source tree has since moved shared numerical/sampling code
into \path{library/spark/}. The 4 November counts remain evidence about
their pinned source snapshot; they are not a GNATprove result for the
refactored bodies. This report revision reuses immutable evidence and
adds exposition and comparisons. Current-source proof requires a new
matching campaign rather than silently substituting historical Ada.


\chapter{Dask execution, reductions and recorded diagnostics}
Dask creates lazy random arrays in chunks of \SimChunk\ observations and
constructs shared tasks for \(G\) and \(E\). Arithmetic, sums, squared
sums, exceedance counts, replicate comparisons and pooled results belong
to the task graph. The graph is computed after its dependent results have
been assembled, following Dask's documented delayed/array approach
\citep{daskdocs,daskArray,daskDelayed}. SciPy supplies distributional
reference calculations inside those tasks. Report tables are written from
the computed records, and XeLaTeX/PGFPlots draws the figures directly from
those tables. There is no direct Matplotlib plotting stage.

The recorded run retains the generator identity, seeds, Dask/NumPy/SciPy
versions, trial counts, chunk layout and input/output hashes. Reproduction
requires those choices: seeded random streams can depend on chunk layout
and library implementation. Seeded reproducibility is a software property;
it is not a mathematical proof that a finite-state generator produces
independent exact normal variables.

For a fixed size, pooled moments use every observation with equal weight.
The reported sample variance uses denominator \(T-1\); combining replicate
means without their within-replicate sums of squares would lose information.
The mean diagnostic divides the observed standardized mean by its exact
standard error \(1/\sqrt T\). Variance diagnostics use the Gaussian
chi-square reference law. \Cref{sec:simulation-statistical-methods}
derives the mean, variance and count reference calculations and explains
their scope. These checks concern the implementation under ideal sampling;
they cannot validate the chosen psychometric assumptions.

The additional Dask latent-prediction diagnostic compares the empirical
mean squared residual with \eqref{eq:hf-mse}, as shown in
\cref{fig:hf-mse}. The next chapter reports the same estimates alongside
the composite moments and raw-average SDs.

\begin{figure}[htbp]
\centering
\begin{tikzpicture}
\begin{loglogaxis}[
 width=.90\linewidth,height=48mm,
 xlabel={Number of tests \(n\)},ylabel={Mean squared factor-prediction error},
 xmin=1,xmax=1000,
 legend style={font=\small,at={(.98,.98)},anchor=north east},grid=major]
\addplot+[blue,mark=*] table[x=n,y=theory_mse,col sep=comma]
 {simulation-data/moments.csv};
\addlegendentry{Exact target}
\addplot+[red,dashed,mark=x] table[x=n,y=mse,col sep=comma]
 {simulation-data/moments.csv};
\addlegendentry{Dask estimate}
\end{loglogaxis}
\end{tikzpicture}
\caption{The model's factor-prediction MSE declines towards zero even
though the raw score mean has a positive limiting population SD.}
\label{fig:hf-mse}
\end{figure}

\chapter{Observed moments and the variance floor}
\begingroup\small

\begingroup\small
\setlength{\tabcolsep}{4pt}
\Needspace{6\baselineskip}
\noindent\textbf{Pooled standardized moments, raw-mean SD and factor-prediction MSE}
\begin{longtable}{rrrrrrr}
\toprule
$n$ & $\bar Z_C$ & $S^2_{Z_C}$ & SD target & SD observed & MSE target & MSE observed \\
\midrule
\endfirsthead
$n$ & $\bar Z_C$ & $S^2_{Z_C}$ & SD target & SD observed & MSE target & MSE observed \\
\midrule
\endhead
\bottomrule
\endfoot
$1$ & $-6.156160\times10^{-4}$ & $1.014525$ & $1.000000$ & $1.007237$ & $0.400000$ & $0.402245$ \\
$2$ & $-0.001659$ & $1.015857$ & $0.894427$ & $0.901491$ & $0.250000$ & $0.251070$ \\
$3$ & $-0.002182$ & $1.016318$ & $0.856349$ & $0.863307$ & $0.181818$ & $0.182512$ \\
$4$ & $-0.002507$ & $1.016523$ & $0.836660$ & $0.843544$ & $0.142857$ & $0.143373$ \\
$5$ & $-0.002733$ & $1.016624$ & $0.824621$ & $0.831447$ & $0.117647$ & $0.118060$ \\
$10$ & $-0.003303$ & $1.016717$ & $0.800000$ & $0.806659$ & $0.062500$ & $0.062713$ \\
$100$ & $-0.004224$ & $1.016275$ & $0.777174$ & $0.783473$ & $0.006623$ & $0.006648$ \\
$1000$ & $-0.004498$ & $1.015972$ & $0.774855$ & $0.781018$ & $6.662225\times10^{-4}$ & $6.689853\times10^{-4}$ \\
\end{longtable}
\endgroup

\endgroup
The standardized mean should fluctuate around zero and its sample variance
around one for every size. In this run, the largest absolute standardized
mean discrepancy is \SimMaxMeanSE\ theoretical standard errors. The
reported table retains the actual deviations, rather than replacing them
with their targets. Differences between neighbouring rows are correlated
because they share random inputs; counting all eight rows as independent
validation would overstate the evidence.

The raw-average SD columns show the contraction in
\eqref{eq:hf-average-variance}. At \(n=1\), raw and standardized scores
coincide. The rapid early decline slows as the \((1-c)/n\) residual
contribution becomes small relative to the common variance \(c\).
At \(n=100\) and \(1000\), adding a further factor of ten in test count
changes the raw population SD only slightly. Figure~\ref{fig:hf-sd}
compares that exact target with the Dask estimate.

\begin{figure}[htbp]
\centering
\begin{tikzpicture}
\begin{semilogxaxis}[
 width=.90\linewidth,height=53mm,
 xlabel={Number of tests \(n\)},ylabel={SD of the raw mean},
 xmin=1,xmax=1000,ymin=.74,ymax=1.02,
 legend style={font=\small,at={(.98,.98)},anchor=north east},
 grid=major]
\addplot+[blue,mark=*] table[x=n,y=theory_sd,col sep=comma]
 {simulation-data/moments.csv};
\addlegendentry{Exact Gaussian target}
\addplot+[red,dashed,mark=x] table[x=n,y=observed_sd,col sep=comma]
 {simulation-data/moments.csv};
\addlegendentry{Dask estimate}
\addplot[black,dotted,domain=1:1000,samples=2] {sqrt(.6)};
\addlegendentry{Limit \(\sqrt{0.6}\)}
\end{semilogxaxis}
\end{tikzpicture}
\caption{Averages retain common population variation. The composite is
standardized separately to SD one; it is not the quantity on this axis.}
\label{fig:hf-sd}
\end{figure}

The mean-squared-error columns and \cref{fig:hf-mse} examine a different prediction target.
For a single test the exact MSE is \(0.4\); at ten tests it is
\(0.0625\); at one thousand it is \(0.4/600.4\). These are direct
substitutions in \eqref{eq:hf-mse}. Their square roots are residual SDs
for predicting \(G\), not population SDs of the score average. Numerical
agreement with these targets demonstrates the distinction in simulated
data; it does not identify a real common factor or show that residual
errors in real tests are independent.

\chapter{Tail counts, rare events and replicate variation}
\begingroup\small

\begingroup\small
\setlength{\tabcolsep}{4pt}
\Needspace{6\baselineskip}
\noindent\textbf{Pooled moderate-tail exceedances}
\begin{longtable}{rrrrr}
\toprule
$n$ & $K_2$ & $K_2/T$ & $K_3$ & $K_3/T$ \\
\midrule
\endfirsthead
$n$ & $K_2$ & $K_2/T$ & $K_3$ & $K_3/T$ \\
\midrule
\endhead
\bottomrule
\endfoot
$1$ & $909698$ & $0.022742$ & $54247$ & $0.001356$ \\
$2$ & $909724$ & $0.022743$ & $53994$ & $0.001350$ \\
$3$ & $909845$ & $0.022746$ & $54113$ & $0.001353$ \\
$4$ & $910036$ & $0.022751$ & $54129$ & $0.001353$ \\
$5$ & $909677$ & $0.022742$ & $54000$ & $0.001350$ \\
$10$ & $909847$ & $0.022746$ & $53908$ & $0.001348$ \\
$100$ & $909953$ & $0.022749$ & $53786$ & $0.001345$ \\
$1000$ & $910060$ & $0.022752$ & $53806$ & $0.001345$ \\
\end{longtable}
\endgroup

\begingroup\small
\setlength{\tabcolsep}{4pt}
\Needspace{6\baselineskip}
\noindent\textbf{Rare-event frequencies and exact bounds, all rates per million}
\begin{longtable}{rrrrr}
\toprule
$n$ & $K_q$ & $10^6 K_q/T$ & Exact 95\% lower & Exact 95\% upper \\
\midrule
\endfirsthead
$n$ & $K_q$ & $10^6 K_q/T$ & Exact 95\% lower & Exact 95\% upper \\
\midrule
\endhead
\bottomrule
\endfoot
$1$ & $4$ & $1.000$ & $0.272$ & $2.560$ \\
$2$ & $3$ & $0.750$ & $0.155$ & $2.192$ \\
$3$ & $3$ & $0.750$ & $0.155$ & $2.192$ \\
$4$ & $3$ & $0.750$ & $0.155$ & $2.192$ \\
$5$ & $3$ & $0.750$ & $0.155$ & $2.192$ \\
$10$ & $1$ & $0.250$ & $0.006$ & $1.393$ \\
$100$ & $3$ & $0.750$ & $0.155$ & $2.192$ \\
$1000$ & $3$ & $0.750$ & $0.155$ & $2.192$ \\
\end{longtable}
\endgroup


\endgroup
% Promote the retained article heading while preserving the recorded data.
\begingroup
\let\subsection\section
\subsection{Specific observations in the recorded run}
At $n=1$ and $n=1000$, the observed factor-prediction MSEs are $0.402245$ and $6.689853\times10^{-4}$, respectively. The raw-mean SD at $n=1000$ is $0.781018$, close to its positive correlation floor $\sqrt{0.6}$; the decreasing MSE concerns recovery of the simulated factor, not removal of population differences. The pooled rare count ranges from 0 to 0 across the eight coupled sizes. A count on this scale is sensitive to individual extreme draws; the exact intervals quantify that uncertainty. The shared draws prevent adding these eight counts as independent evidence.

\endgroup
The three exceedance columns are nested for each draw:
\(\mathbf1_{Z>q}\leq\mathbf1_{Z>3}\leq\mathbf1_{Z>2}\), since
\(q>3>2\). Their counts must obey the same ordering. The probabilities
in \eqref{eq:hf-tail-targets} are constant across sizes, so systematic
changes in the raw-average cutoff are compatible with essentially the
same standardized exceedance rates. The table compares empirical rates
with the appropriate fixed distribution, rather than applying raw-score
cutoffs to already standardized scores.

At the two moderate thresholds, the expected counts in the pooled sample
are \(T\overline\Phi(2)\) and \(T\overline\Phi(3)\).
Their binomial standard errors are small relative to those counts.
The one-in-a-million threshold has only \(T\times10^{-6}=4\) expected
exceedances per fixed size, with count SD
\(\sqrt{4(1-10^{-6})}\), approximately two. The observed pooled rare
counts range from \SimRareMin\ to \SimRareMax\ across the eight sizes.
That range is reported as observed; it is not eight independent estimates
because the comparisons share \(G,E\).

A rare count of zero does not show that the true probability is zero.
Under the model a single one-million-observation replicate has zero rare
exceedances with probability \((1-10^{-6})^{10^6}\), approximately
\(e^{-1}\). The pooled four-million-observation probability is
\((1-10^{-6})^{4\times10^6}\), approximately \(e^{-4}\).
The exact probability follows by multiplying the independent failure
probabilities; the exponential approximation follows from
\(\log(1-p)=-p+O(p^2)\). Confidence bounds in the table preserve
uncertainty when counts are very small.
Proposition~\ref{hfpr:simulation} proves both the exact zero-count law
and a non-asymptotic bound on the stated exponential approximations.

\begingroup\small

\begingroup\small
\setlength{\tabcolsep}{4pt}
\Needspace{6\baselineskip}
\noindent\textbf{Individual seeded runs; maxima are over the eight coupled sizes}
\begin{longtable}{rrrrr}
\toprule
Replicate & Seed & Max $|z_{mean}|$ & Max $|S^2-1|$ & $K_q$ at $n=1000$ \\
\midrule
\endfirsthead
Replicate & Seed & Max $|z_{mean}|$ & Max $|S^2-1|$ & $K_q$ at $n=1000$ \\
\midrule
\endhead
\bottomrule
\endfoot
$1$ & $20261114$ & $1.190$ & $7.389498\times10^{-4}$ & $6$ \\
$2$ & $20261115$ & $0.380$ & $5.374933\times10^{-4}$ & $7$ \\
$3$ & $20261116$ & $0.622$ & $0.001204$ & $3$ \\
$4$ & $20261117$ & $1.192$ & $9.342210\times10^{-4}$ & $5$ \\
$5$ & $20261118$ & $0.747$ & $4.131702\times10^{-4}$ & $9$ \\
$6$ & $20261119$ & $0.614$ & $0.001205$ & $4$ \\
$7$ & $20261120$ & $1.288$ & $4.263499\times10^{-4}$ & $5$ \\
$8$ & $20261121$ & $0.248$ & $6.967057\times10^{-4}$ & $4$ \\
\end{longtable}
\endgroup

\endgroup
Replicate-level output shows information that pooling can hide: the
particular seed's mean/variance fluctuation, the presence or absence of
rare events, and the variation of counts between runs. The replicate
index denotes a fresh seeded stream, while battery size denotes a
coupled transformation within that replicate. Results are not selected
from seeds that happen to look closest to the model.

For Bernoulli exceedances, the empirical frequency has variance
\(p(1-p)/T\), because it is an average of independent zero-one
indicators. Its relative standard error is therefore
\(\sqrt{(1-p)/(Tp)}\). At \(p=10^{-6}\), the present expected count
of four gives roughly \(50\%\) relative standard error. Even a
\(10\%\) relative \emph{standard error}, rather than a 95\% interval
half-width, would require roughly one hundred million observations
per fixed size. This requirement follows by solving
\(\sqrt{(1-p)/(Tp)}\leq0.1\) for \(T\). More importantly, increasing
simulation size reduces Monte Carlo error only within the assumed model;
it supplies no empirical evidence that a human-score tail is Gaussian.
The exact precision inequality and count-moment proof are in
Proposition~\ref{hfpr:simulation} and
Proposition~\ref{sm:binomial}; interval coverage is addressed in
Results~\ref{sm:wilson} and \ref{sm:clopper-pearson}.

\chapter{Comparing the SPARK and Dask experiments}
\label{sec:hf-dual-comparison}
Both implementations simulate \eqref{eq:hf-aggregate} and standardize with
the same \(v_n\). Both report the unbiased sample variance with denominator
\(T-1\). Thus their composite means, variances and exceedance frequencies
have the same ideal targets, although their finite random streams and
reduction orders differ. The mathematical equivalence follows from
Proposition~\ref{hfpr:factor}; the source-pinned driver and replay records
establish which formulas these particular native executions used.

The Dask pool has four times the observations of the native experiment
for each fixed size. By Proposition~\ref{sm:normal-sample}, its ideal
mean SE is \(0.0005\), compared with \(0.001\) in SPARK; the variance SE
is \(\sqrt{2/(T-1)}\), approximately \(0.000707\) and \(0.001414\),
respectively. The variance-SE formula follows directly from that
proposition's chi-square law and \(\operatorname{Var}(\chi^2_\nu)=2\nu\).
This predicts greater precision at the larger sample size, rather than
superiority of a programming language or generator.

At \(n=1\), SPARK records mean \(0.000396961\) and variance
\(0.999022835\); Dask records mean \(0.000188666\) and variance
\(1.000254209\). At \(n=1000\), the respective means are
\(0.000729021\) and \(0.000330421\), and variances are
\(1.001649339\) and \(1.000209530\). Every native and pooled Dask mean
is less than one of its own theoretical SEs from zero. These are
descriptions of the two recorded samples, not additional independent
experiments at each plotted size. Different seeds and generators alone
would not justify treating cross-implementation differences as independent
normal errors; no such difference test is used here.

At \(n=1\), native moderate-tail counts are \(22603\) and \(1312\),
versus pooled Dask counts \(90719\) and \(5544\); their denominators are
one million and four million. The resulting rates fluctuate on either
side of the same targets. A useful caution is the Dask \(Z>3\) interval
for this size: its pointwise 95\% Wilson lower endpoint is approximately
\(0.001350017\), just above the target \(0.001349898\). The table retains
this discrepancy. A nominal pointwise interval is not guaranteed to cover
on every run, and multiple coupled diagnostics are not a simultaneous
95\% guarantee; see Propositions~\ref{sm:wilson} and \ref{sm:multiple}.

The native rare counts are zero or one, while pooled Dask counts range
from one to four. Their expected counts are one and four. From
Proposition~\ref{sm:binomial}, rare-frequency relative SE is
\(\sqrt{(1-p)/(Tp)}\): approximately 100\% for SPARK and 50\% for the
Dask pool at \(p=10^{-6}\). This proves why agreement of tiny counts is
weak evidence about the numerical experiment's tail and supplies no
empirical validation of a psychometric extreme tail.

The native CSV does not record latent-factor squared errors. Dask's MSE
columns therefore test an additional model consequence, rather than
replicate an unreported SPARK measurement. SPARK's deterministic
calculation output does check the corresponding conditional-factor
formulas and worked examples, as described in
\cref{sec:hf-spark-calculations}. Neither implementation observes real
latent ability. Runtime tests, native arithmetic proofs, limited Python
verification and statistical-model assumptions keep their separate scopes.

\chapter{Interpretation and practical limits}
The computations illustrate a specific mathematical distinction.
Standardization keeps composite population variance at one. Correlated
component averages retain a common source of variation. A conditional
predictor can nevertheless become more accurate for a latent variable
when the model assumes that additional residual components are genuinely
independent. These conclusions follow from the stated model and the
proved variance/regression calculations; simulation supplies a numerical
reproduction of them.

They do not establish the common-factor structure, the population
correlation, valid score comparability, practice-effect independence, or
a one-in-a-million frequency in real test takers. The testing standards
treat population calibration and the justification of score
interpretations as evidential obligations \citep[Standards 1.1--1.4,
1.14 and 5.8--5.9]{standards2014}. The present simulation contains no
observations capable of discharging those obligations. Its detailed
counts and uncertainty are useful precisely because they identify what
has, and has not, been demonstrated.

\appendix
\chapter{Referenced foundations and full working proofs}
\label{sec:proof-foundations}

The results below supply the mathematical foundations used in this report.
The explicitly assumed inputs are finite second moments and, where stated,
a jointly Gaussian law and independent observations. An assumption is not a
consequence of a simulation. Definitions, source quotations, historical
records, rounded numerical evaluations and recommendations have the status
specified in the claim map; they are not additional mathematical theorems.
The proofs reduce the deductions to linearity of integration, the spectral
theorem, the Gaussian characteristic function and conditional expectation.
Source locators refer to the linked author versions specified in the
bibliography, rather than silently mixing their pagination with print editions.

\noindent\textbf{Document-revision provenance.}
The current book-format revision is recorded in
\path{reporting/book-format/revision.json}, with current TeX/Bib sources
and PDF hashes. The earlier proof-text revision remains recorded in
\path{reporting/proof-revision/revision.json}.
The original \path{reporting/run.json} retains the numerical
simulation campaign, its original document hashes and numerical outputs.
The book revision checks and reuses the unchanged simulation-data files; it does
not rerun the experiment or claim a new native SPARK/Rocq proof.
The repository's subsequent shared-library refactor is separate from the
retained native campaigns. Their source-pinned results do not constitute
new GNATprove evidence for refactored Ada bodies; those changes require
their own matching campaign. The written mathematical proofs below do not
depend on that program-verification status.

\begin{proposition}[Finite moments and affine changes of scale]
\label{pf:moments}
Let a real random vector \(X\) have finite second moments, mean \(\mu\)
and covariance \(\Sigma\). Let \(Y\) be square-integrable, \(w\) a
deterministic vector of matching dimension, and \(b,s\) real constants. Then
\[
 \E(w^{\mathsf T}X)=w^{\mathsf T}\mu,\qquad
 \Var(w^{\mathsf T}X)=w^{\mathsf T}\Sigma w,\qquad
 \SD(b+sY)=|s|\SD(Y).
\]
A zero variance is equivalent to equality with the mean almost surely.
\end{proposition}
\begin{proof}
H\"older's inequality ensures integrability of all products
\((X_i-\mu_i)(X_j-\mu_j)\). Finite linearity of the integral then gives
the expectation formula and
\[
 \E\!\left[\Big(\sum_iw_i(X_i-\mu_i)\Big)^2\right]
 =\sum_{i,j}w_iw_j\E[(X_i-\mu_i)(X_j-\mu_j)]
 =w^{\mathsf T}\Sigma w.
\]
Replacing \(Y\) by \(b+sY\) multiplies its centered value by \(s\), so
its variance is \(s^2\Var(Y)\); taking the nonnegative square root proves
the SD formula. If \(\E[(Y-\E Y)^2]=0\), Markov's inequality applied to
\((Y-\E Y)^2\) makes \(P(|Y-\E Y|>\epsilon)=0\) for every
\(\epsilon>0\). The countable union over \(\epsilon=1/n\) proves the
almost-sure assertion. The reverse implication follows by integration.
These are applications of linearity, H\"older and the variance identities
in \citet[Section 1.6, especially Theorems 1.6.1, 1.6.3--1.6.4]{durrett2019}.
\end{proof}

\begin{proposition}[Covariance feasibility and Gaussian construction]
\label{pf:spectral}
A real symmetric matrix is a covariance matrix if and only if it is
positive semidefinite. For the equicorrelation matrix
\(R_n=(1-c)I_n+c\mathbf1_n\mathbf1_n^{\mathsf T}\), \(n\geq2\), this
is equivalent to \(-1/(n-1)\leq c\leq1\). Positive definiteness requires
both inequalities to be strict.
\end{proposition}
\begin{proof}
Proposition~\ref{pf:moments} gives
\(w^{\mathsf T}\Sigma w=\E[(w^{\mathsf T}(X-\mu))^2]\geq0\).
Conversely, the real spectral theorem diagonalizes a symmetric positive
semidefinite matrix as \(QDQ^{\mathsf T}\), with orthogonal \(Q\) and
nonnegative diagonal entries \(D\). For independent standard normal
coordinates \(Z\), the vector \(QD^{1/2}Z\) has covariance
\(QD^{1/2}I D^{1/2}Q^{\mathsf T}=\Sigma\) and is Gaussian, including
zero eigenvalues. The spectral theorem and positive square roots are
\citet[Theorems 7.29 and 7.38]{axler2024}; the Gaussian construction is
also \citet[Theorems 3.4.4--3.4.5]{gallager2013}.
For \(R_n\), the span of \(\mathbf1_n\) has eigenvalue
\(1+(n-1)c\), and its orthogonal complement has eigenvalue \(1-c\),
because \(\mathbf1_n^{\mathsf T}w=0\) there. Every vector decomposes
into these two orthogonal subspaces. Nonnegativity, or positivity, of
both eigenvalues gives precisely the stated inequalities. If one fixed
\(c\) must work for arbitrarily large \(n\), the lower bounds force
\(c\geq0\); the finite-dimensional bounds alone permit negative \(c\).
\end{proof}

\begin{proposition}[Gaussian closure and block independence]
\label{pf:gaussian}
An affine image of a jointly Gaussian vector is jointly Gaussian. Its law
is determined by its mean and covariance. Two blocks of a jointly Gaussian
vector are independent whenever their cross-covariance is zero, including
singular Gaussian vectors.
\end{proposition}
\begin{proof}
For a Gaussian vector with mean \(\mu\) and covariance \(\Sigma\), its
characteristic function is
\(\psi_X(t)=\exp(i t^{\mathsf T}\mu-t^{\mathsf T}\Sigma t/2)\).
Thus \(\psi_{AX+b}(t)=e^{it^{\mathsf T}b}\psi_X(A^{\mathsf T}t)\)
has the Gaussian form with mean \(A\mu+b\) and covariance
\(A\Sigma A^{\mathsf T}\). This also follows from the independent-normal
affine representation, without invertibility. If a two-block covariance
has zero off-diagonal blocks, the displayed characteristic function
factors as \(\psi_U(s)\psi_V(t)\). By Fubini this is the characteristic
function of the product of the two marginal measures. Uniqueness of
characteristic functions proves both independence and determination by
mean and covariance. The Gaussian formula and closure are
\citet[Theorems 3.3.1--3.3.2, equations (3.14) and (3.18)]{gallager2013};
uniqueness is \citet[Theorem 3.10.4]{durrett2019}.
\end{proof}

\begin{proposition}[Gaussian conditional distribution, including exact histories]
\label{pf:conditioning}
Suppose \((U,V)\) is jointly Gaussian with mean \((\mu_U,\mu_V)\),
covariance blocks \(A,B\) and \(C=\Cov(U,V)\), and \(B\) positive
definite. A version of the conditional law, for every \(v\), is Gaussian
with mean and covariance
\[
 \mu_U+CB^{-1}(v-\mu_V),\qquad A-CB^{-1}C^{\mathsf T}.
\]
The conditional-mean formula is the unique continuous version on the
whole conditioning space. The full joint covariance need not be invertible.
\end{proposition}
\begin{proof}
Set \(L=CB^{-1}\) and
\(W=U-\mu_U-L(V-\mu_V)\). Its cross-covariance with \(V\) is
\(C-LB=0\). Expanding its covariance gives
\(A-LC^{\mathsf T}-CL^{\mathsf T}+LBL^{\mathsf T}
=A-CB^{-1}C^{\mathsf T}\). Proposition~\ref{pf:gaussian} makes
\(W\) Gaussian and independent of \(V\). Consequently the kernel
\(K(v,H)=P(\mu_U+L(v-\mu_V)+W\in H)\) satisfies
\[
 P(U\in H,V\in D)=\int_D K(v,H)P_V(dv)
\]
by independence and Fubini, which is the defining conditional-law identity.
Its expectation is the displayed affine function. A positive-definite
Gaussian \(V\) has a strictly positive density everywhere. If two
continuous versions differed at any point, continuity would give a
neighborhood of positive probability on which they differ, contradicting
almost-sure uniqueness. Conditioning on an exact continuous history uses
this specified version, not division by the probability-zero event
\(\{V=v\}\). This residual proof is the Gaussian regression theorem
\citep[Theorems 3.5.1--3.5.2]{gallager2013}, proved here in the slightly
more general form allowing singular \((U,V)\); conditional versions are
defined in \citet[Section 4.1]{durrett2019}.
\end{proof}

\begin{proposition}[Normal marginal distributions do not imply joint normality]
\label{pf:normal-marginals}
Normal marginals and a specified Pearson correlation do not in general
determine Gaussian composite or joint-tail laws.
\end{proposition}
\begin{proof}
Let \(X\sim N(0,1)\) and independently let \(S\) equal \(1\) with
probability \((1+\rho)/2\) and \(-1\) otherwise, \(-1<\rho<1\).
Put \(Y=SX\). Symmetry gives \(Y\sim N(0,1)\), and independence gives
\(\E[XY]=\E[S]\E[X^2]=\rho\). Nevertheless
\(P(X+Y=0)=(1-\rho)/2>0\). A nondegenerate normal variable has a
density and no atoms; \(X+Y\) has variance \(2+2\rho>0\), so this
pair is not jointly Gaussian. For a concrete tail nonidentification at
\(\rho=0\), let \(t>0\) and \(p=P(X>t)\in(0,1/2)\). This pair
has \(P(X>t,Y>t)=p/2\), whereas an independent standard-normal pair,
with the same marginals and zero correlation, has intersection \(p^2\).
These differ since \(p<1/2\). The construction uses symmetry,
independence and the product-integration rule
\citep[Sections 1.7 and 2.1]{durrett2019}; it is a counterexample,
not an alternative assumption secretly used for the report's Gaussian tables.
\end{proof}

\begin{proposition}[Event arithmetic and conditional denominators]
\label{pf:events}
For events \(A,B\) with probabilities \(p_A,p_B\) and intersection
\(J\), their union has probability \(p_A+p_B-J\). If its conditioning
probability is positive, \(P(B\mid A)=J/p_A\). The correlation of the
nonconstant indicators is
\[
 \frac{J-p_Ap_B}{\sqrt{p_A(1-p_A)p_B(1-p_B)}}.
\]
\end{proposition}
\begin{proof}
The pointwise identities
\(1_{A\cup B}=1_A+1_B-1_{A\cap B}\) and
\(1_A1_B=1_{A\cap B}\) prove the union and covariance formulas by
linearity. Because \(1_A^2=1_A\),
\(\Var(1_A)=p_A-p_A^2\), and similarly for \(B\); divide their
covariance by the product of their positive SDs. The conditional formula
is precisely the event-conditional probability definition, requiring
\(p_A>0\). The four disjoint cells have probabilities
\(J,p_A-J,p_B-J,1-p_A-p_B+J\), obtained by the same identities.
These are consequences of finite additivity and conditional probability
\citep[Sections 1.1 and 4.1]{durrett2019}.
\end{proof}

\begin{proposition}[Tower property and population--conditional variance]
\label{pf:expectation}
For square-integrable \(Y\) and a sub-\(\sigma\)-algebra \(\mathcal G\),
\[
 \E\E[Y\mid\mathcal G]=\E Y,\qquad
 \Var(Y)=\E\Var(Y\mid\mathcal G)
              +\Var(\E[Y\mid\mathcal G]).
\]
\end{proposition}
\begin{proof}
The defining conditional-integral identity with the whole space gives
the first formula. Write \(M=\E[Y\mid\mathcal G]\). Conditional
Jensen \citep[Theorem 4.1.10]{durrett2019} ensures \(M\) is
square-integrable. In
\(Y-\E Y=(Y-M)+(M-\E Y)\), the cross term has expectation zero:
the pull-out and tower properties give
\(\E[(Y-M)(M-\E Y)]
=\E[(M-\E Y)\E(Y-M\mid\mathcal G)]=0\).
Expanding the square and integrating gives the variance identity. The
pull-out rule extends from bounded multipliers by truncation and
Cauchy--Schwarz. These conditional-expectation properties follow from
\citet[Section 4.1, especially Theorems 4.1.13--4.1.14]{durrett2019}.
In particular, fixing an observed starting score and varying it across
the population are different probability experiments, even when both
use the same marginal test SD.
\end{proof}

\section{Proofs of the composite and prediction claims}
\label{hfpr:appendix}

This appendix proves the mathematical implications used in this report.
The hypotheses are part of each result: a proof under a Gaussian model
does not prove that model for human scores. The foundations are linearity
of expectation and conditional expectation, covariance expansion and the
Gaussian characteristic function \citep[Sections 1.6, 3.3 and 4.1]{durrett2019},
Gaussian affine closure and regression \citep[Theorems 3.3.1--3.3.2 and
3.5.1--3.5.2]{gallager2013}, and the spectral theorem for real symmetric
matrices \citep[Theorem 7.29]{axler2024}. Their relevant forms and proofs
are stated in the shared mathematical foundations preceding this appendix.
In particular, Propositions~\ref{pf:moments}, \ref{pf:spectral},
\ref{pf:gaussian}, \ref{pf:conditioning} and \ref{pf:expectation}
state the moment, spectral, Gaussian and conditional-expectation tools.
Write \(\phi(x)=(2\pi)^{-1/2}\exp(-x^2/2)\),
\(\Phi(x)=\int_{-\infty}^x\phi(t)\,dt\), and
\(\overline\Phi=1-\Phi\).

\begin{proposition}[Covariance, scales and weighted composites]
\label{hfpr:standardization}
For square-integrable scores \(X_i\) with means \(\mu_i\) and positive
SDs \(\sigma_i\), put \(Z_i=(X_i-\mu_i)/\sigma_i\). Then
\(\E Z_i=0\), \(\Var(Z_i)=1\), and their covariance matrix is their
correlation matrix \(R\). For any fixed real vector \(w\),
\[
 \E(w^{\mathsf T}Z)=0,\qquad
 \Var(w^{\mathsf T}Z)=w^{\mathsf T}Rw.
\]
When the latter is positive, division by its square root gives a mean-zero,
unit-variance composite. The transformation \(I=100+sZ_C\), \(s>0\),
has population mean 100 and SD \(s\).
\end{proposition}
\begin{proof}
Linearity gives \(\E(X_i-\mu_i)=0\). Expanding the square in the
definition of variance gives \(\Var(aX+b)=a^2\Var(X)\); this proves
the first two conclusions and, on expanding a product, gives
\(\Cov(Z_i,Z_j)=\Cov(X_i,X_j)/(\sigma_i\sigma_j)=r_{ij}\).
For any finite collection, integrability of each product follows from
Cauchy--Schwarz. Hence finite sums can be taken through expectation:
\[
 \Var\!\left(\sum_iw_iZ_i\right)
 =\E\!\left[\sum_{i,j}w_iw_jZ_iZ_j\right]
 =\sum_{i,j}w_iw_jr_{ij}=w^{\mathsf T}Rw.
\]
For two unstandardized scores the same centered expansion gives
\(\Var(X+Y)=\Var(X)+\Var(Y)+2\Cov(X,Y)\).
The assertions about division and reporting scales now follow by the
same affine-variance identity. No distributional assumption was used.
For example, \((160-100)/15=4\) and \((160-100)/16=15/4=3.75\).
\end{proof}

\begin{proposition}[Equal weights, admissibility and information in the matrix]
\label{hfpr:equal-weights}
For \(n\geq2\), let \(\bar r=2\sum_{i<j}r_{ij}/[n(n-1)]\). Then
\[
 v_n:=\Var(\bar Z)=\frac{1+(n-1)\bar r}{n}
 =\bar r+\frac{1-\bar r}{n},\qquad
 -\frac1{n-1}\leq\bar r\leq1.
\]
When \(v_n>0\), the composite is \(\bar Z/\sqrt{v_n}\).
The entire correlation matrix must be positive semidefinite (PSD);
the displayed bound on its average is necessary, rather than a sufficient
test for an arbitrary proposed matrix. For an equicorrelation matrix the
same bound is necessary and sufficient for PSD. The average correlation
alone need not determine the variance with unequal weights.
\end{proposition}
\begin{proof}
Apply Proposition~\ref{hfpr:standardization} with every weight \(1/n\).
The diagonal contributes \(n/n^2\); symmetry makes the off-diagonal
contribution \(2\sum_{i<j}r_{ij}/n^2\). Substitution proves the formula.
Every quadratic form \(u^{\mathsf T}Ru=\Var(u^{\mathsf T}Z)\) is
nonnegative, proving PSD and, with \(u\) the all-ones vector, the lower
bound on \(\bar r\). Cauchy--Schwarz gives \(|r_{ij}|\leq1\), proving
the upper bound. For \(n=1\), directly \(v_1=1\).

In the equicorrelated case \(R=(1-c)I+c\mathbf1\mathbf1^{\mathsf T}\).
Its action on \(\mathbf1\) multiplies by \(1+(n-1)c\); its action
on every vector orthogonal to \(\mathbf1\) multiplies by \(1-c\).
The spectral theorem therefore gives PSD exactly for
\(-1/(n-1)\leq c\leq1\). Such a PSD matrix is a realizable Gaussian
covariance: diagonalize it, take nonnegative square roots of its
eigenvalues, and apply the resulting square-root matrix to independent
standard normals. A zero \(v_n\) makes the mean deterministic and the
division undefined.

To prove the unequal-weight qualification, compare \(R_0=I_3\) with
\[
 R_t=\begin{pmatrix}1&t&-t\\t&1&0\\-t&0&1\end{pmatrix},\qquad t=1/4.
\]
Both average off-diagonal correlations are zero. The latter matrix has
eigenvalues \(1,1+\sqrt2t,1-\sqrt2t\), all positive, as is checked
by its characteristic polynomial
\((1-\lambda)[(1-\lambda)^2-2t^2]\).
For \(w=(1,1,0)^{\mathsf T}\), however,
\(w^{\mathsf T}R_tw=2+2t\ne2=w^{\mathsf T}R_0w\).
Finally, if \(I_i=100+sZ_i\), then
\(\bar Z=(\bar I-100)/s\). Substitution into
\(100+s\bar Z/\sqrt{v_n}\) proves the IQ-style formula.
\end{proof}

\begin{proposition}[Gaussian percentiles require a joint model]
\label{hfpr:normality}
If \(Z\) is jointly Gaussian and \(w^{\mathsf T}Rw>0\), its
standardized fixed weighted composite is standard normal and its
model-based percentile is \(100\Phi(z_C)\). Neither two moments nor
normal marginal distributions imply this conclusion.
\end{proposition}
\begin{proof}
An affine image of a jointly Gaussian vector is Gaussian
\citep[Theorem 3.3.1]{gallager2013}; Proposition~\ref{hfpr:standardization}
gives its mean zero and variance one. By definition of its CDF, the fraction
at or below \(z_C\) is then \(\Phi(z_C)\).

For the failure of the converse, take \(X\sim N(0,1)\) and an independent
sign \(S\) with \(P(S=1)=(1+c)/2\) and \(P(S=-1)=(1-c)/2\),
where \(-1<c<1\). Set \(Y=SX\). Symmetry of the normal distribution
makes \(Y\) standard normal. Independence gives
\(\Cov(X,Y)=\E(SX^2)=\E S\,\E X^2=c\).
Nevertheless \(P(X+Y=0)=(1-c)/2>0\): on \(S=-1\) the sum vanishes,
whereas on \(S=1\) it is \(2X\), which has no atom. A nondegenerate
normal distribution has a density and hence no atoms. Standardization
preserves this atom, so these normal marginals and their correlation do
not imply a normal composite. They also fail to identify joint tail
probabilities: for positive \(a,b\), this model gives
\(P(X>a,Y>b)=(1+c)\overline\Phi(\max(a,b))/2\), since opposite signs
cannot both exceed positive cutoffs. It is not the Gaussian joint law.
The same counterexample, having the required moments, proves that moments
alone do not determine the tail.
\end{proof}

\begin{proposition}[Thresholds, monotonicity and fixed-correlation limits]
\label{hfpr:thresholds}
Let \(q=\Phi^{-1}(1-10^{-6})\) be defined by the continuous standard
normal law. For a Gaussian composite with mean-variance factor \(v_n>0\),
the one-in-a-million threshold on the component mean is
\(t_n(c)=q\sqrt{c+(1-c)/n}\). For \(0\leq c\leq1\) it decreases
with \(n\), increases with \(c\) for \(n>1\), and tends to
\(q\sqrt c\). For fixed positive observed mean, its standardized
composite decreases with \(c\). These are properties of the stipulated
sequence of batteries, rather than a bound on possible score values.
\end{proposition}
\begin{proof}
Since \(\phi>0\), integration shows that \(\Phi\) is continuous and
strictly increasing from 0 to 1. The quantile is uniquely defined and
\(P(Z_C>q)=1-\Phi(q)=10^{-6}\). Because \(\sqrt{v_n}>0\),
\(\bar Z>q\sqrt{v_n}\) is exactly \(Z_C>q\), proving the cutoff.
The normal law has no atoms, so strict or non-strict inequalities have
the same model probability.

The ordering \(q>3>2\), used for nested simulated counts, can be shown
without relying on its printed decimal. On \([3,4]\) the density is
at least \(\phi(4)\), and
\[
 \overline\Phi(3)\geq\int_3^4\phi(x)\,dx\geq\phi(4)
 =\frac{1}{\sqrt{2\pi}e^8}>
 \frac{1}{3\cdot3^8}=\frac1{19683}>10^{-6}.
\]
Here \(e=\sum_{j=0}^{\infty}1/j!<3\), by comparison of terms
\(j!\geq2^{j-1}\) for \(j\geq1\) with strict inequality for
\(j\geq3\), and \(\pi=4\int_0^1(1+x^2)^{-1}\,dx<4\), so
\(\sqrt{2\pi}<3\). Strict decrease of the normal upper tail proves
\(q>3\).

For real \(x\geq1\) put \(v(x,c)=c+(1-c)/x\). Elementary
differentiation gives
\[
 \frac{\partial t}{\partial x}
 =-\frac{q(1-c)}{2x^2\sqrt{v(x,c)}},\qquad
 \frac{\partial t}{\partial c}
 =\frac{q(1-1/x)}{2\sqrt{v(x,c)}}.
\]
The signs prove the claims, with equality in the evident singleton or
perfect-correlation cases. For a fixed \(b>0\),
\(\partial[b/\sqrt{v(x,c)}]/\partial c
=-b(1-1/x)/(2v(x,c)^{3/2})\leq0\).
For \(c<1\), the raw-SD derivative has magnitude
\((1-c)/(2\sqrt{x^3(cx+1-c)})\). The positive denominator increases
with \(x\geq1\) for \(0\leq c<1\), proving the slowing decline;
for \(c=1\) it is constant.
Also \(v_n-c=(1-c)/n\to0\); continuity of square root gives the limit.
When \(c>0\), rationalizing yields the explicit error bound
\[
 0\leq t_n(c)-q\sqrt c
 =\frac{q(1-c)}{n(\sqrt{v_n}+\sqrt c)}
 \leq\frac{q(1-c)}{2n\sqrt c}.
\]
For \(c=0\) the exact cutoff is \(q/\sqrt n\to0\). Fixed negative
\(c\) cannot give valid equicorrelated batteries for arbitrarily large
\(n\), by Proposition~\ref{hfpr:equal-weights}. Under the Gaussian model
every nondegenerate \(Z_C\) has unbounded support; thus a limit on the
component-mean \emph{cutoff} is not a ceiling on attainable composites.
Changing the correlations as tests are added changes the sequence and
does not satisfy the fixed-\(c\) hypothesis. Finally \(100+s q\) is
the composite reporting-scale cutoff for every \(n\), whereas
\(100+s t_n(c)\) is the component-average cutoff, by the affine identity.
\end{proof}

\begin{proposition}[Common-factor aggregate, posterior and residual contrasts]
\label{hfpr:factor}
Assume \(0<c<1\) and
\(X_i=\sqrt c\,G+\sqrt{1-c}\,\varepsilon_i\), where all displayed
basic variables are independent standard normals. For
\(M_n=n^{-1}\sum_iX_i\), put \(v_n=c+(1-c)/n\) and
\(R_n=c/v_n=nc/[1+(n-1)c]\). Then
\[
 (G,M_n)\ \stackrel d=\
 \left(G,\sqrt c\,G+\sqrt{(1-c)/n}\,E\right),
\]
where \(E\) is standard normal independent of \(G\). Given the entire
score vector, \(G\) is normal with mean
\(\widehat G_n=(\sqrt c/v_n)M_n=\sqrt{R_n}\,Z_{C,n}\) and variance
\(1-R_n=(1-c)/[1+(n-1)c]\). Moreover
\[
 \operatorname{Corr}(G,Z_{C,n})^2=R_n,\quad
 \Var(\widehat G_n)=R_n,\quad
 \E[(G-\widehat G_n)^2]=1-R_n.
\]
The conditional mean minimizes mean-square prediction error among all
square-integrable functions of the scores.
\end{proposition}
\begin{proof}
Linearity, independence and unit variances give \(\E X_i=0\),
\(\Var(X_i)=1\), \(\Cov(X_i,X_j)=c\) for \(i\ne j\), and
\(\Cov(G,M_n)=\sqrt c\). Define
\(E=n^{-1/2}\sum_i\varepsilon_i\). Its Gaussian affine law has mean
zero and variance one, and it is independent of \(G\) because it is a
function of residuals independent of \(G\). Substitution gives the
stated representation, including its joint law with \(G\).

Let \(W=G-(\sqrt c/v_n)M_n\). Its variance, obtained by expanding
the square, is \(1+c/v_n-2c/v_n=1-c/v_n\), and
\(\Cov(W,M_n)=\sqrt c-(\sqrt c/v_n)v_n=0\). For each contrast
\(D_i=X_i-M_n\), symmetry gives
\[
 \Cov(G,D_i)=\sqrt c-\sqrt c=0,\qquad
 \Cov(M_n,D_i)=\frac{1+(n-1)c}{n}-v_n=0.
\]
Thus \(\Cov(W,D_i)=0\) as well. All these variables are jointly
Gaussian. Their characteristic function is
\(\exp(it^{\mathsf T}\mu-t^{\mathsf T}\Sigma t/2)\); zero
cross-covariance makes it factor between \(W\) and \((M_n,D_1,\ldots,D_n)\).
Uniqueness of characteristic functions proves independence even though
the contrasts are linearly dependent. This is the Gaussian regression
principle \citep[Theorems 3.3.2 and 3.5.1--3.5.2]{gallager2013}.
Since the score vector is recovered by \(X_i=M_n+D_i\), \(W\) is
independent of all observed scores. Conditional on them,
\(G=(\sqrt c/v_n)M_n+W\), with the stated normal mean and variance.
This explicitly proves the claim that the contrasts add no information
about \(G\) under this symmetric Gaussian model.

The covariance with \(Z_{C,n}=M_n/\sqrt{v_n}\) is
\(\sqrt{c/v_n}\). Both \(G\) and \(Z_{C,n}\) have variance one,
proving the squared-correlation formula. Scaling \(M_n\) gives
\(\Var(\widehat G_n)=c/v_n\). The centered independent residual gives
the MSE formula. For any square-integrable score function \(h\),
conditional expectation and pull-out give
\(\E[W(\widehat G_n-h)]=0\), since \(\E[W\mid X]=0\).
Expanding \(G-h=W+(\widehat G_n-h)\) proves
\[
 \E[(G-h)^2]=1-R_n+\E[(\widehat G_n-h)^2]\geq1-R_n.
\]
Finally \(1-R_n=(1-c)/(nc+1-c)\to0\) for fixed \(c>0\), although
\(\SD(M_n)=\sqrt{v_n}\to\sqrt c>0\). Population spread of scores
and conditional prediction error are thus different mathematical objects.
\end{proof}

\begin{proposition}[Observed score laws do not identify a latent interpretation]
\label{hfpr:nonidentification}
Even an exactly known observable correlation matrix, or the entire joint
score law, need not identify a proposed latent ability, its posterior, or
an individual measurement-error interval.
\end{proposition}
\begin{proof}
Keep precisely the same observable \(X\) from
Proposition~\ref{hfpr:factor}. In one model call its common factor \(G\)
the target. In a second model extend the probability space by a standard
normal \(H\) independent of \((G,X)\), and call \(H\) the target.
Every observable distribution and its matrix \(R\) are unchanged, yet
\(\E[H\mid X]=0\), \(\Var(H\mid X)=1\), whereas the first model has
the different posterior proved in Proposition~\ref{hfpr:factor}.
Therefore no rule depending only on those observable quantities can
deduce both latent posteriors correctly. Additional measurement assumptions
or external validity evidence are logically required. The example does not
claim that either target is the correct psychometric construct.
\end{proof}

\begin{proposition}[Selection and repeated-score procedures can alter calibration]
\label{hfpr:selection}
Population correlations need not equal selected-sample correlations.
Likewise the maximum or a stopping rule based on repeated attempts need
not have the distribution of one prespecified score. The direction and
size of an adjustment require assumptions about selection and dependence.
\end{proposition}
\begin{proof}
Let independent \(X,E\) be standard normal and let
\(Y=\rho X+\sqrt{1-\rho^2}\,E\), \(0<\rho<1\).
Unconditionally \(\operatorname{Corr}(X,Y)=\rho\).
Select on \(X>0\). Symmetry gives \(\E[X^2\mid X>0]=1\), while
\[
 \E[X\mid X>0]=2\int_0^\infty x\phi(x)\,dx
 =2\phi(0)=\sqrt{2/\pi},
\]
using \(\phi'(x)=-x\phi(x)\). Consequently
\(v=\Var(X\mid X>0)=1-2/\pi\in(0,1)\).
Independence of \(E\) from the selection event gives conditional
covariance \(\rho v\) and variance of \(Y\) equal to
\(\rho^2v+1-\rho^2\). The selected correlation is therefore
\[
 \rho\sqrt{\frac{v}{\rho^2v+1-\rho^2}}<\rho,
\]
because \(\rho^2v+1-\rho^2-v=(1-\rho^2)(1-v)>0\).
No selection gives the original value, so there is already no universal
correction independent of the mechanism. This constructive mathematical
example supports the possibility of distortion; empirical correction
methods are attributed to \citet{sackett2000}.

For independent standard-normal attempts \(U_1,\ldots,U_m\),
independence gives
\(P(\max_jU_j\leq t)=\Phi(t)^m\), so the chance of any qualification
above a fixed cutoff is \(1-\Phi(t)^m\), exceeding the one-attempt chance
for finite \(t\) and \(m>1\). The probability of no success after \(m\)
attempts tends to zero, proving eventual qualification with probability
one if attempts can continue independently without limit. In contrast,
perfectly repeated scores \(U_j=U_1\) give qualification probability
\(1-\Phi(t)\) for every \(m\). These are explicit counterexamples to
treating all procedures as interchangeable; no claim that every possible
dependent procedure differs is needed. Practice effects themselves are
an empirical finding attributed to \citet{hausknecht2007}.

Finally a finite number \(m\) of observations cannot logically determine
an unknown rare-event probability: the event of no exceedances has
probability \((1-p)^m>0\) for every \(0<p<1\) in an independent
Bernoulli model. At \(m=27\) and \(p=10^{-6}\), the union bound gives
\(P(\hbox{no event})\geq1-27\cdot10^{-6}=0.999973\).
Selection further changes the sampling experiment. A small selected
sample is therefore not a deductive establishment of a one-in-a-million
population frequency.
\end{proof}

\begin{proposition}[Coupling, rare-event precision and an exponential bound]
\label{hfpr:simulation}
The aggregate simulation preserves each fixed-size \((G,M_n)\) law but
reusing a residual across sizes does not simulate a nested battery.
Nested exceedance indicators imply nested counts. Under independent
ideal draws with event probability \(p>0\), relative Monte Carlo standard
error is \(\sqrt{(1-p)/(Tp)}\). For \(0<p<1\),
\[
 e^{-Tp-Tp^2/[2(1-p)]}\leq(1-p)^T\leq e^{-Tp}.
\]
\end{proposition}
\begin{proof}
The aggregate representation is proved in Proposition~\ref{hfpr:factor}.
With one shared \(E\), two sizes \(m<n\) have covariance
\(c+(1-c)/\sqrt{mn}\). In a genuine nested battery their residual
averages share only the first \(m\) terms, giving instead
\(c+(1-c)/n\), by covariance expansion. These differ for \(0<c<1\),
so the fixed-size shortcut is not a claim about nested paths. As a concrete
dependence diagnostic, two aggregates formed with the shared \((G,E)\)
have positive covariance; they cannot be independent centered variables.
The per-size laws still follow from the separately proved representation.

For \(q>3>2\), every real \(z\) satisfies
\(\mathbf1_{z>q}\leq\mathbf1_{z>3}\leq\mathbf1_{z>2}\).
Summing pointwise proves the event-count ordering without any sampling
assumption. For independent Bernoulli indicators, covariance expansion
gives \(\Var(K/T)=p(1-p)/T\); divide its SD by \(p\) to obtain the
relative error. For \(p=10^{-6}\) and \(T=4\cdot10^6\) it is
\(\sqrt{(1-10^{-6})/4}\), approximately one half. A target relative
SE at most \(\delta\) requires exactly
\(T\geq(1-p)/(p\delta^2)\); at \(\delta=0.1\) this is
\(T\geq99{,}999{,}900\). This is a standard-error calculation, not a
95\% confidence-width guarantee.

Independence gives \(P(K=0)=(1-p)^T\). The fundamental theorem of
calculus yields
\[
 -\log(1-p)-p=\int_0^p\frac{t}{1-t}\,dt,
 \qquad0\leq\int_0^p\frac{t}{1-t}\,dt
 \leq\frac{p^2}{2(1-p)}.
\]
Multiply by \(T\), exponentiate and preserve the inequalities to obtain
the stated bound. For \(p=10^{-6}\) and \(T=10^6\) or \(4\cdot10^6\),
the exponent correction is at most \(5.000006\cdot10^{-7}\) or
\(2.000003\cdot10^{-6}\), respectively; hence the quoted
\(e^{-1}\) and \(e^{-4}\) approximations have explicitly controlled
relative error. Binomial count moments, Gaussian moment diagnostics,
pooled variance and coverage of the reported intervals are proved in the
shared statistical-reference appendix (Propositions~\ref{sm:normal-sample},
\ref{sm:binomial} and \ref{sm:pooling}, and the interval results
\ref{sm:sd-interval}, \ref{sm:wilson}, \ref{sm:clopper-pearson}); those laws assume ideal independent
draws, not merely a seeded finite-state generator.
For the empirical factor MSE specifically, the independent residuals across
draws obey \(W_j/\sqrt{1-R_n}\sim N(0,1)\), by
Proposition~\ref{hfpr:factor}. Therefore
\(\sum_{j=1}^T W_j^2/(1-R_n)\) is chi-square with \(T\) degrees of
freedom by its defining sum-of-squares law. This also identifies its
expectation \(T\) and thus the empirical MSE target \(1-R_n\).
\end{proof}

\begin{proposition}[Exact arithmetic behind the illustrative numbers]
\label{hfpr:numerics}
For three SD-16 scores averaging 160 and \(c=3/5\),
\[
 v_3=\frac{11}{15},\qquad
 Z_C=\sqrt{\frac{3375}{176}},\qquad
 I_C=100+\sqrt{\frac{54000}{11}}.
\]
For the latent example when \(Z_C=q\) at the exact quantile,
\[
 R_3=\frac9{11},\qquad
 \E[G\mid X]=\frac3{\sqrt{11}}q,\qquad
 \SD(G\mid X)=\sqrt{\frac2{11}}.
\]
The threshold table consists exactly of \(q\sqrt{(3n+2)/(5n)}\)
and its affine SD-16 transform; its infinite row is \(q\sqrt{3/5}\).
\end{proposition}
\begin{proof}
Substitute \(\bar z=15/4\), \(c=3/5\), \(n=3\) into
Propositions~\ref{hfpr:equal-weights} and \ref{hfpr:factor}. Squaring
the positive worked composite gives
\((15/4)^2/(11/15)=3375/176\); squaring its SD-16 excess gives
\(16^2\cdot3375/176=54000/11\). Exact rational comparisons give
\[
 (4.37905)^2<3375/176<(4.37915)^2,\qquad
 (70.055)^2<54000/11<(70.065)^2.
\]
Monotonicity of the positive square root proves rounding to \(4.3791\)
and \(170.06\). For the posterior,
\(R_3=(9/5)/(11/5)=9/11\) and \(1-R_3=2/11\), proving the formulas.
The rational inequalities
\((0.425)^2<2/11<(0.435)^2\) and
\((6.815)^2<512/11<(6.825)^2\) prove the displayed conditional-SD
roundings \(0.43\) and \(6.82\).
The scale conversion is \(100+16(3/\sqrt{11})q\) for its mean and
\(16\sqrt{2/11}\) for its SD. The SD-15 singleton threshold is
\(100+15q\); the three-test threshold at \(c=9/10\) is
\(100+16q\sqrt{14/15}\). These exact expressions account for every
worked decimal and threshold-table entry.

The displayed quantile and its resulting decimal values are numerical
approximations, reproduced by the documented R upper-tail quantile or
\texttt{scipy.stats.norm.isf} and the repository's calculation script.
Agreement of these software evaluations is numerical evidence, not by
itself an exact-real proof of their final rounded digits. Sensitivity to
quantile error is explicit: replacing \(q\) by \(q+\eta\) changes a
component threshold by exactly \(\eta\sqrt{v_n}\), its SD-\(s\) version
by \(s\eta\sqrt{v_n}\), and the latent-scale mean by
\(16\sqrt{R_n}\eta\). Thus the algebra is proved independently of
backend accuracy, and any certified quantile enclosure can be propagated
by these positive factors. Likewise recorded sample means, variances,
counts and interval endpoints are observations/computed approximations,
not identities asserted of all realizations.
\end{proof}

\chapter{Statistical reference laws and uncertainty calculations}
\label{sec:shared-statistical-methods}
\label{sec:simulation-statistical-methods}

These results concern ideal independent observations from the stated model.
They justify reference laws and interval calculations, rather than proving
that finite-state pseudorandom output is independent or that the model fits
real tests. The input theorems are Gaussian affine closure and the Gaussian
characteristic function \citep[Theorems 3.3.1--3.3.2]{gallager2013},
orthonormal-basis extension \citep[Theorem 6.36]{axler2024}, and linearity,
independence and elementary probability inequalities
\citep[Sections 1.6, 2.1 and 3.4]{durrett2019}. The derivations below do
not depend on a successful software test.

\begin{proposition}[Normal sample moments and the chi-square pivot]
\label{sm:normal-sample}
For an integer \(m\geq2\), independent \(Y_i\sim N(\mu,v)\), \(v>0\),
the sample mean is \(N(\mu,v/m)\). With
\(S^2=\sum_i(Y_i-\bar Y)^2/(m-1)\),
\((m-1)S^2/v\sim\chi^2_{m-1}\), independently of \(\bar Y\), and
\(\E S^2=v\).
\end{proposition}
\begin{proof}
Finite linearity and zero cross-covariances give the mean and variance;
Gaussian affine closure gives the law. Set \(Z_i=(Y_i-\mu)/\sqrt v\),
and extend the unit vector \(m^{-1/2}(1,\ldots,1)\) to an orthonormal
basis. Let \(Q\) have this vector as its first row. The transformed
Gaussian \(T=QZ\) has covariance \(QQ^{\mathsf T}=I\); its characteristic
function factors, so all its coordinates are independent standard normals.
Its first coordinate is \(\sqrt m(\bar Y-\mu)/\sqrt v\). Orthogonality
preserves squared norm, hence
\[
 \frac{1}{v}\sum_i(Y_i-\bar Y)^2
 =\sum_iZ_i^2-T_1^2=\sum_{j=2}^mT_j^2.
\]
The right side is chi-square by definition, independent of \(T_1\).
Its expectation is \(m-1\), proving unbiasedness of \(S^2\).
This argument supplies the normal and chi-square laws directly from the
referenced Gaussian and linear-algebra theorems.
\end{proof}

\begin{corollary}[SD intervals and fixed-parameter reference bands]
\label{sm:sd-interval}
Let \(q_u\) denote a \(\chi^2_{m-1}\) quantile and \(0<\alpha<1\).
An exact \(1-\alpha\) confidence interval for the unknown SD under this
normal model is
\[
 \left[\sqrt{\frac{(m-1)S^2}{q_{1-\alpha/2}}},
       \sqrt{\frac{(m-1)S^2}{q_{\alpha/2}}}\right].
\]
For a fixed theoretical \(v\), a reference interval for \(S^2\) is instead
\([vq_{\alpha/2}/(m-1),vq_{1-\alpha/2}/(m-1)]\). The reference
mean discrepancy in standard-error units is
\((\bar Y-\mu)/\sqrt{v/m}\sim N(0,1)\).
\end{corollary}
\begin{proof}
The continuous chi-square law gives
\(P(q_{\alpha/2}\leq(m-1)S^2/v\leq q_{1-\alpha/2})=1-\alpha\).
Solving these inequalities for \(v\), then taking the increasing positive
square root, gives the SD interval. Solving for \(S^2\), keeping \(v\)
fixed, gives the reference interval. The last assertion is the standardized
sample-mean law in Proposition~\ref{sm:normal-sample}. If \(v=0\), every
observation equals \(\mu\) almost surely, so this division-based pivot is
inapplicable; constant output is checked directly. SciPy evaluates the
quantiles \citep{scipystats}; the algebra, rather than the library call,
establishes the probability statement.
\end{proof}

\begin{proposition}[Bernoulli counts and sampling precision]
\label{sm:binomial}
For an integer \(m\geq1\) and \(m\) independent event indicators of common
probability \(p\in[0,1]\), the
count \(K\) is binomial, \(\E K=mp\), and
\(\Var(K/m)=p(1-p)/m\).
\end{proposition}
\begin{proof}
Independence makes the probability generating function of their sum
\((1-p+pt)^m\). The binomial theorem gives
\(P(K=j)=\binom mjp^j(1-p)^{m-j}\). Each indicator has mean \(p\)
and variance \(p(1-p)\), because its square equals itself. Summing means
and independent variances, then dividing by \(m^2\), proves the claims.
These are finite-product independence and linearity
\citep[Sections 1.6 and 2.1]{durrett2019}.
\end{proof}

\begin{proposition}[Wilson score endpoints and their coverage scope]
\label{sm:wilson}
For an integer \(m\geq1\), an observed count \(K\in\{0,\ldots,m\}\),
\(0<\alpha<1\), \(\hat p=K/m\) and
\(z=\Phi^{-1}(1-\alpha/2)\), inversion
of the normal score inequality gives endpoints
\[
 \frac{\hat p+z^2/(2m)\ \mathord\pm\
 z\sqrt{\hat p(1-\hat p)/m+z^2/(4m^2)}}{1+z^2/m}.
\]
The endpoint algebra is exact; nominal \(1-\alpha\) coverage uses a normal
score approximation and is not an exact finite-sample coverage theorem.
\end{proposition}
\begin{proof}
Squaring \(|\hat p-p|\leq z\sqrt{p(1-p)/m}\) gives
\[
 (m+z^2)p^2-(2m\hat p+z^2)p+m\hat p^2\leq0.
\]
The positive leading coefficient means precisely the interval between its
two real roots. Its discriminant is
\(4mz^2\hat p(1-\hat p)+z^4\). The quadratic formula, dividing
numerator and denominator by \(2m\), gives the stated endpoints. For
\(0<p<1\), the classical independent-identically-distributed central
limit theorem implies
\(\sqrt m(\hat p-p)/\sqrt{p(1-p)}\Rightarrow N(0,1)\), giving limiting
coverage at its continuity points, not equality for fixed \(m\).
This is the score construction of \citet{wilson1927}; the limit theorem is
\citet[Theorem 3.4.1]{durrett2019}. The implementation reuses
SciPy's binomial-result interval routine \citep{scipyBinomtest}.
\end{proof}

\begin{proposition}[Exact equal-tail binomial intervals]
\label{sm:clopper-pearson}
For an integer \(m\geq1\), a count \(j\in\{0,\ldots,m\}\), and
\(0<\alpha<1\), define \(L(0)=0\), \(U(m)=1\); for \(j>0\),
let \(L(j)\) solve \(P_{L(j)}(K\geq j)=\alpha/2\), and for \(j<m\),
let \(U(j)\) solve \(P_{U(j)}(K\leq j)=\alpha/2\). Then
\(P_p(L(K)\leq p\leq U(K))\geq1-\alpha\) for every \(p\in[0,1]\).
\end{proposition}
\begin{proof}
For \(0<j\leq m\), differentiating the finite binomial sum and
telescoping gives
\[
 \frac{d}{dp}P_p(K\geq j)
 =m\binom{m-1}{j-1}p^{j-1}(1-p)^{m-j}>0\quad(0<p<1).
\]
The tail runs continuously from zero to one, so \(L(j)\) exists uniquely;
the lower tail is strictly decreasing and gives \(U(j)\) similarly.
For \(0<j<m\), at \(p=L(j)\) the lower tail is at least
\(1-\alpha/2>\alpha/2\); its decrease implies \(U(j)>L(j)\).
The endpoint counts are ordered directly, so these really are intervals.
For a fixed \(p\), \(L(j)>p\) implies \(P_p(K\geq j)<\alpha/2\).
Such \(j\)'s form an upper terminal set, because upper tails decrease
with \(j\). If nonempty its probability is the tail at its smallest
member and is at most \(\alpha/2\). The analogous lower terminal set
\(U(j)<p\) has probability at most \(\alpha/2\). The union bound makes
total noncoverage at most \(\alpha\); endpoint parameters follow directly.
Discreteness explains why the inequality can be strict. This is the
test-inversion construction of \citet{clopper1934}, with its coverage
proof supplied here. Equivalently the endpoints are beta quantiles:
the derivative above is the \(\mathrm{Beta}(j,m-j+1)\) density and its
tail has value zero at \(p=0\), so integration gives the same cumulative
distribution; the complementary tail gives the upper endpoint.
\end{proof}

\begin{proposition}[Conditional frequencies with random denominators]
\label{sm:conditional-count}
For independent repetitions of a four-cell event experiment, condition on
\(N_A=\ell>0\) observations in event \(A\), where \(P(A)>0\).
The count also in \(B\) is then \(\mathrm{Binomial}(\ell,P(B\mid A))\).
Thus a binomial interval conditional on its observed denominator has its
stated exact or approximate coverage scope. A zero denominator gives no
defined conditional-frequency estimate.
\end{proposition}
\begin{proof}
Let \(q=P(A)\) and \(r=P(B\mid A)\). Grouping cells into
\(A\cap B,A\setminus B,A^c\), their probabilities are
\(qr,q(1-r),1-q\). Product independence and the multinomial counting
coefficient give
\[
 P(N_{AB}=j,N_A=\ell)
 =\binom m\ell\binom\ell j q^\ell(1-q)^{m-\ell}
           r^j(1-r)^{\ell-j}.
\]
Divide by \(P(N_A=\ell)=\binom m\ell q^\ell(1-q)^{m-\ell}\)
when positive. The remaining factor is the stated binomial mass. Cases
with \(q=1\) follow directly with \(\ell=m\). Conditional coverage for
each positive \(\ell\) also gives coverage averaged over positive
denominators by total probability. No convention assigning a number to
\(0/0\) supplies data when \(\ell=0\).
\end{proof}

\begin{proposition}[Pooling counts and moments]
\label{sm:pooling}
For a nonempty finite collection of groups of integer sizes \(m_j\geq2\),
means \(\bar Y_j\) and sample variances
\(S_j^2\), let \(M=\sum_jm_j\) and
\(\bar Y=M^{-1}\sum_jm_j\bar Y_j\). The pooled sample variance satisfies
\[
 (M-1)S^2=\sum_j(m_j-1)S_j^2+\sum_jm_j(\bar Y_j-\bar Y)^2.
\]
Independent equal-design Bernoulli groups pool to a binomial count of size
\(M\); independence is required for that probability law, not for the
displayed algebraic identity.
\end{proposition}
\begin{proof}
Within group \(j\), write
\(Y_{ji}-\bar Y=(Y_{ji}-\bar Y_j)+(\bar Y_j-\bar Y)\).
Squaring and summing, its cross term vanishes because
\(\sum_i(Y_{ji}-\bar Y_j)=0\). Summing groups gives the identity.
For counts, independence multiplies their generating functions to
\(\prod_j(1-p+pt)^{m_j}=(1-p+pt)^M\), proving the binomial law by
Proposition~\ref{sm:binomial}. Sharing random numbers across cases does not
establish this independence premise; in the overlapping-factor designs used
here the rows cannot be counted as additional independent replications.
Averaging group SDs generally does
not satisfy the displayed within-plus-between identity.
\end{proof}

\begin{proposition}[Simultaneous diagnostics and rare-event precision]
\label{sm:multiple}
If \(J\) fixed diagnostics have false-alarm probabilities at most
\(\alpha_j\), the probability of any false alarm is at most
\(\sum_{j=1}^J\alpha_j\), regardless of their dependence.
\end{proposition}
\begin{proof}
The pointwise bound \(1_{\cup_j F_j}\leq\sum_j1_{F_j}\), integrated,
is the union bound. Taking each \(\alpha_j=\alpha/J\) therefore gives
a family bound \(\alpha\). The same argument applies to two tails per
row, counting both in \(J\). Finite diagnostics remain observations
under a reference model and cannot establish an entire distribution.
The bound is the elementary subadditivity theorem
\citep[Section 1.1]{durrett2019}.
\end{proof}

\begin{corollary}[Zero counts, relative precision and sample-size scaling]
\label{sm:rare}
For independent Bernoulli trials, \(P(K=0)=(1-p)^m\). For \(p>0\) the
relative standard error of \(K/m\) is \(\sqrt{(1-p)/(mp)}\); requiring
it at most \(\epsilon>0\) means \(m\geq(1-p)/(p\epsilon^2)\).
At fixed \(p\in(0,1)\), multiplying \(m\) by four halves the standard
error. Zero successes do not prove \(p=0\).
\end{corollary}
\begin{proof}
The first statement is the binomial mass at zero. Divide the SD in
Proposition~\ref{sm:binomial} by \(p\), square the nonnegative inequality,
and rearrange to get the sample-size bound. The square-root dependence
proves the scaling. Every \(0<p<1\) assigns positive probability to
zero successes for finite \(m\), proving the final statement. A desired
95\% zero-count upper confidence bound, for example, follows from
\((1-U)^m=.025\) for the equal-tail construction, not from setting
\(p\) equal to zero. These conclusions concern sampling uncertainty;
they do not bound floating-point quadrature error or model error.
\end{proof}

\chapter{Claim-to-proof and evidence map}
\label{hfmap:simulation}

The report separates exact deductions under ideal sampling, computed
approximations, observations in recorded files, software provenance and
external empirical claims. In particular, ``exact target'' means an exact
mathematical expression under the model; its printed decimal and SciPy
evaluation still involve numerical approximation. The following inventory
covers every model, diagnostic and interpretation used in the report.

\begingroup\small
\setlength{\tabcolsep}{4pt}
\begin{longtable}{>{\raggedright\arraybackslash}p{.35\linewidth}>{\raggedright\arraybackslash}p{.59\linewidth}}
\toprule
Claim/location & Proof or evidence and precise scope\\
\midrule
\endfirsthead
\toprule Claim/location & Proof or evidence and precise scope\\\midrule
\endhead
\bottomrule\endfoot
Experimental units, \(c=0.6\), battery sizes and independent ideal base
normals &
Explicit simulation design/model hypotheses. They are not inferred from
the historical selected scorers or justified by passing software tests.\\
Individual score means, SDs and correlations;
equations~\ref{eq:hf-factor-model}--\ref{eq:hf-average-variance} &
Propositions~\ref{hfpr:standardization}, \ref{hfpr:equal-weights} and
\ref{hfpr:factor} prove every covariance and the finite-size variance.\\
Aggregate shortcut, joint law with the latent factor and its limits as a
representation; equation~\ref{eq:hf-aggregate} &
Proposition~\ref{hfpr:factor} proves the joint law by summing independent
normal residuals. Proposition~\ref{hfpr:simulation} proves the distinction
between shared-residual coupling and genuinely nested battery averages.\\
Standardization, normality and tail targets;
equations~\ref{eq:hf-composite}--\ref{eq:hf-tail-targets} &
Propositions~\ref{hfpr:standardization}, \ref{hfpr:normality} and
\ref{hfpr:thresholds}. The strict normal CDF defines the exact quantile;
SciPy computes approximate reference probabilities/quantiles.\\
Raw cutoffs, reporting scales, decreasing SD and cutoff, positive SD
floor and unit composite SD &
Propositions~\ref{hfpr:thresholds} and \ref{hfpr:standardization}, by
event equivalence, affine variance and the fixed-\(c\) limit. Unbounded
support proves that a threshold floor does not cap possible scores.\\
Predictor, reliability-like coefficient, posterior SD and MSE;
equations~\ref{eq:hf-predictor}--\ref{eq:hf-mse} &
Proposition~\ref{hfpr:factor}, including independent Gaussian residuals,
conditioning on all contrasts, MSE optimality, predictor variance and
the zero-MSE limit. The model's factor is specified, not discovered.\\
MSEs at \(n=1,10,1000\), raw and standardized scores coincide at
\(n=1\), contraction slows &
Substitution gives \(2/5\), \(1/16\) and \(1/1501\), respectively,
from Proposition~\ref{hfpr:factor}. Proposition~\ref{hfpr:thresholds}
proves the decreasing derivative magnitude and positive floor; direct
\(v_1=1\) proves the singleton identity.\\
Dask construction, chunks, seeds, versions, actual reductions and figures &
Recorded software design and source/output hashes in
\texttt{reporting/run.json}, the shared reporting code and generated
\texttt{simulation-data} files; Dask's APIs are documented in
\citet{daskdocs,daskArray,daskDelayed}. Reproducibility is conditional on
the recorded implementation choices, not a proof of exact normal RNGs.\\
Pooled sample variance and mean/SD diagnostics &
Propositions~\ref{sm:pooling} and \ref{sm:normal-sample}, and
Corollary~\ref{sm:sd-interval}, prove the reductions and reference laws.
Proposition~\ref{hfpr:simulation} gives the MSE sum-of-squares law.\\
Observed tables, chart points and discussion of largest discrepancies,
variance errors, SDs and MSEs &
Observations/computed summaries in \texttt{simulation-data/moments.csv},
\texttt{replicate-records.csv} and their generated TeX files. Their exact
targets are proved above; the realized values are not universal
theorems. Statistical compatibility is an inference under the ideal
sampling laws, with no guarantee that every diagnostic passes.\\
Nested moderate/rare event counts and constant standardized tail targets &
The pointwise indicator ordering is proved in
Proposition~\ref{hfpr:simulation}; the Gaussian tail equality is
Proposition~\ref{hfpr:thresholds}. Ordering requires no independence.\\
Expected counts, count SDs, rate SEs, relative precision and required trial
count; zero-event probabilities and exponential approximations &
Propositions~\ref{sm:binomial} and \ref{hfpr:simulation}, and
Corollary~\ref{sm:rare}; the simulation proof derives an explicit exponential error
bound and the exact \(99{,}999{,}900\) trial lower bound for 10\% relative
SE at \(p=10^{-6}\).\\
Wilson intervals for moderate events; exact binomial intervals for rare
events &
Propositions~\ref{sm:wilson} and \ref{sm:clopper-pearson} prove score
inversion and conservative exact-test coverage. The Wilson coverage is
approximate; the exact interval's coverage theorem does not certify
rounding of its computed endpoints.\\
Eight coupled sizes are not eight independent validation runs; pooled
replicates and multiple diagnostics &
Proposition~\ref{hfpr:simulation} demonstrates dependence of the score
arrays. Proposition~\ref{sm:pooling} applies across ideal independent
replicates, and Proposition~\ref{sm:multiple} gives a union-bound
interpretation of several diagnostics without assuming independent rows.\\
Python/Dask composite simulation, seeded replicates and coupled sizes &
The recorded \path{simulation-data/moments.csv} and
\path{simulation-data/tails.csv} retain the per-size sample moments
and exceedance frequencies. Proposition~\ref{hfpr:simulation}
explains why shared draws across sizes are not independent.
The generator is not assumed to establish a human population model.\\
Model-based sampling precision and rare-event uncertainty &
Propositions~\ref{sm:normal-sample}, \ref{sm:binomial},
\ref{sm:wilson}, \ref{sm:clopper-pearson} and
\ref{sm:multiple} establish the ideal Gaussian/binomial
precision and diagnostic-coverage statements. Computed finite
intervals may miss their nominal target; many coupled diagnostic
intervals are not a simultaneous 95\% confidence statement.\\
Python/Deal contracts and mathematical proof boundaries &
The Dask source, simulation configuration and retained numerical
tables document the implementation. Deal lint and tests check
specified Python properties; Rocq checks precisely stated formal
theorems. No archived compiler run, third-party implementation proof
or direct psychometric data validation is claimed.\\
Lack of real-model validation, latent-factor identification and
population-tail evidence &
Propositions~\ref{hfpr:normality} and \ref{hfpr:nonidentification}
give counterexamples to identification from moments or score laws;
\ref{hfpr:selection} proves selection/procedure counterexamples.
The need for empirical validity evidence is attributed to
\citet[Standards 1.1--1.4, 1.14, 5.8--5.9]{standards2014}.\\
AI disclosure and no participant data &
Process/provenance statements, not mathematical deductions or claims
of independent human review.\\
\end{longtable}
\endgroup

\begingroup
\setlength{\bibsep}{3pt}
\bibliographystyle{plainnat}
\bibliography{simulation_references,references,hf_simulation_references}
\endgroup
\clearpage
\chapter{AI disclosure and provenance}\label{sec:ai-disclosure}
% Compatibility alias: canonical AI-use disclosure is in library/latex.
% Canonical full appendix text; the chapter heading is defined in report_style.tex.
This report and its supporting repository were substantially generated and
revised with OpenAI Codex and collaborating AI agents. The session identifies
Codex as based on the GPT-6 family. Exact serving checkpoints, per-agent model
identifiers and generation settings were not preserved; no more precise model
attribution is claimed. This disclosure covers the mathematical exposition,
simulations, programs, documentation, verification, typesetting and publication,
including subsequent layout and documentation revisions.

AI contributions included interpreting the problem and historical articles;
developing and reviewing derivations, counterexamples, theorem references,
proofs and claim maps; calculating examples and limits; editing
Python/Dask/Deal implementations, contracts, tests and auditors; executing
the recorded checks; producing tables and scientific charts; preparing the
XeLaTeX/KOMA-Script books; visually reviewing PDFs; and preparing repository
updates.

The requester Anders Hellström supplied questions and sources, clarified notation and
averaging, selected example parameters, specified formats and verification
requirements, requested revisions and authorized publication. Repository
ownership and commit attribution do not identify every passage as human-authored
or independently verified. Parallel mathematical and visual reviews were
performed by AI agents and automated tools within the same workflow.
Independent human peer review, professional sign-off and endorsement by
original authors, source hosts or cited organizations are not recorded and
are not claimed.

Original authors and references remain credited. Historical articles are
analysed rather than relicensed; replacement articles are AI-generated
analytical reconstructions. Sources were consulted at different levels,
including mathematical texts, articles, official software documentation and
bibliographic metadata. Reference notes identify versions and access limits;
full-text examination of every cited work is not claimed. Mathematical proofs
establish consequences of stated hypotheses, rather than empirical adequacy
of the IQ model.

Verification claims rest on retained commands, actual outputs, tool versions,
source hashes and checksums. Previous disconnected implementations and
their proof campaigns were removed from maintained source; they remain accessible
through earlier Git history and are not current-source verification for
Python/Dask. Deal contract declarations, static lint, runtime/property
tests, exact Rocq kernel checks, numerical accuracy, PRNG properties and
model validity have separate scopes. Unsupported, skipped or unrun checks
are not reported as successful proofs. Trusted numerical/library boundaries
and specification adequacy remain explicit.
Independent numerical validation means a separate implementation, algorithm or
library, not independent human authorship or review.

No participant data were collected and no empirical psychometric validation
study was conducted.
Simulations use synthetic observations under stated probability models. A typesetting-only revision does not, by itself, execute additional random
simulations or prove additional Rocq theorems. Older
compiler evidence is available through Git history, not active source files.
The repository does not preserve a complete prompt/tool transcript, exact per-line authorship
or AI-contribution percentage; these have not been invented. Its license
applies to original material, subject to third-party rights.


The 16 November 2026 revision used AI agents to expand the SPARK study,
compare the native and Dask results, audit retained observations and review
the rebuilt pages. No new random experiment or native prover run was performed.
\end{document}
